\pdfoutput=1

\documentclass[a4paper,twoside,leqno]{article}

% \documentclass[a4paper,draft,twoside,leqno]{article}
% \documentclass[a4paper,twoside,leqno]{article}

\usepackage{a4wide}

\usepackage[p,osf]{scholax}
% T1 and textcomp are loaded by package. Change that here, if you want
% load sans and typewriter packages here, if needed
\usepackage{amsmath,amsthm}% must be loaded before newtxmath
% amssymb should not be loaded
\usepackage[scaled=1.075,ncf,vvarbb]{newtxmath}% need to scale up math package
% vvarbb selects the STIX version of blackboard bold.


% The usual ams packages
% \usepackage{amssymb}
% \usepackage{amsmath}
% \usepackage{mathtools}          % for multilined environment
% \usepackage{amsthm}

% % for incrementing two counters at the same time
% \usepackage{xassoccnt}

% for the \vvvert command
% \usepackage{mathabx}

% microtype for better typesetting
\usepackage{microtype}

% for \llbracket and \rrbracket
\usepackage[only,llbracket,rrbracket]{stmaryrd}

% To produce nice calligraphic letters for Lebesgue and Hausdorff measure with
% '\mathscr{L}'
% \usepackage{mathrsfs} %% !!!!

% To insert figures.
\usepackage{graphicx}

% for mean value integral sign
% \usepackage{esint}                % load after the AMS packages

% for different enumeration styles
\usepackage{enumitem}
% \setlist[enumerate,1]{label=(\alph*),resume}
\setlist[enumerate,1]{label=(\alph*)}
\setlist[enumerate,2]{label=(\roman*)}

\usepackage[T1]{fontenc}
\usepackage[utf8x]{inputenc}
\usepackage[usenames,dvipsnames]{xcolor}

\usepackage{url}
\usepackage[hidelinks,colorlinks=true,linkcolor=blue,citecolor=darkgray,linktocpage=false]{hyperref}

% To show only referenced labels
\usepackage{mathtools}
\mathtoolsset{showonlyrefs}

% For the symbol \mathds{1} used for the characteristic function
\usepackage{dsfont}

% \usepackage{phonetic}
% \usepackage{extpfeil}
% \usepackage[all]{xy}

% to ease referencing
% \usepackage[notref,notcite]{showkeys}

% No hyphenation is better than an odd looking one
\hyphenation{Lip-schitz-ian Besicovitch}

% "empty" theorem enviroment for use in the introduction
\newtheoremstyle{citing}% name
  {3pt}%      Space above, empty = `usual value'
  {3pt}%      Space below
  {\itshape}% Body font
  {}%         Indent amount (empty = no indent, \parindent = para indent)
  {\bfseries}% Thm head font
  {.}%        Punctuation after thm head
  {.5em}%     Space after thm head: " " = normal interword space;
        %       \newline = linebreak
  {\thmnote{#3}}% Thm head spec
\theoremstyle{citing}
\newtheorem*{citing}{}

% theorem, lemma, corollary, remark, example and definition as usual
% also, intro_problem to describe the problem dealt with in the introduction
% and miniremark as environment without name to gather assumptions and the
% like
\theoremstyle{definition}
\newtheorem*{intro_problem}{Problem}

\newcounter{introcnt}
\newcounter{claimcount}
\theoremstyle{remark}
\newtheorem{claim}[claimcount]{Claim}
\newtheorem{subclaim}{Claim}[claimcount]

\swapnumbers
\theoremstyle{plain}

\newtheorem{introthm}[introcnt]{Theorem}
\newtheorem{introcol}[introcnt]{Corollary}

\newtheorem{theorem}{Theorem}[section]
\newtheorem{conjecture}[theorem]{Conjecture}
\newtheorem{lemma}[theorem]{Lemma}
\newtheorem{proposition}[theorem]{Proposition}
\newtheorem{corollary}[theorem]{Corollary}

\theoremstyle{remark}
\newtheorem{remark}[theorem]{Remark}
\newtheorem*{remark*}{Remark}
\newtheorem{example}[theorem]{Example}

\theoremstyle{definition}
\newtheorem{definition}[theorem]{Definition}
\newtheorem*{definition*}{Definition}
\newtheorem{miniremark}[theorem]{}
\newtheorem{question}[theorem]{Question}
\newtheorem*{question*}{Question}

% If equation numbers are used at all they are Roman as the enumeration is
% arabic
% \renewcommand{\theequation}{\Roman{equation}}
%\numberwithin{equation}{theorem}

% some counters
\newcounter{counter1}
\newcounter{counter2}

% Commands to conveniently subdevide proofs into several steps
\theoremstyle{definition}
\swapnumbers
\newtheorem{step}{Step} % the environment for steps
\newenvironment{proofinsteps}{\begin{proof}\setcounter{step}{0}}{\end{proof}}

% to reset equation counter each time the theorem counter increases
\usepackage{chngcntr}

% % for \lefthalfcup and \righthalfcup
% \usepackage{MnSymbol}

% \usepackage{lmodern}

% To have the same numbering for equations and enumerates
%
% Warning: putting a numbered equation inside myenumerate may spoil the
% numbering!
%
\newenvironment{myenumerate}{
  \begin{enumerate}
      \setcounter{enumi}{\value{equation}}
    }{
      \setcounter{equation}{\value{enumi}}
  \end{enumerate}
}

% Equation numbers should reset on each new item (theorem, remark, lemma, etc.).
\counterwithin*{equation}{theorem}
\counterwithin*{enumi}{theorem}


% Local Variables:
% coding: utf-8
% eval: (custom-set-variables  '(LaTeX-indent-level 4) '(LaTeX-item-indent -4) '(indent-tabs-mode nil))
% eval: (ispell-change-dictionary "british")
% eval: (flyspell-mode)
% End:


% \newtheorem{introtheorem}{Theorem}
% \renewcommand{\theintrotheorem}{\Alph{introtheorem}} % "letter-numbered" theorems

% \swapnumbers
% \theoremstyle{plain}
% \newtheorem{theorem}{Theorem}[section]
% \newtheorem{lemma}[theorem]{Lemma}
% \newtheorem{claim}[theorem]{Claim}
% \newtheorem{corollary}[theorem]{Corollary}

% \theoremstyle{remark}
% \newtheorem{remark}[theorem]{Remark}
% \newtheorem*{remark*}{Remark}
% \newtheorem{example}[theorem]{Example}

% \theoremstyle{definition}
% \newtheorem{definition}[theorem]{Definition}
% \newtheorem{miniremark}[theorem]{}

% Hom(R^n,R^n)
% \newcommand{\End}[1]{ \mathcal{Z}_{#1} }
% \newcommand{\Endd}[1]{ \mathcal{Z}^*_{#1} }
\newcommand{\End}[1]{ \mathrm{End}({#1}) }
% \newcommand{\EndR}[1]{ \End{\R^{#1}} }
\newcommand{\Endd}[1]{ \mathrm{End}({#1})^* }
% \newcommand{\EndRd}[1]{ \Endd{\R^{#1}} }
% \newcommand{\End}[1]{ \mathrm{End}(#1) }
% \newcommand{\Endd}[1]{ \mathrm{End}(#1)^* }

\newcommand{\SymEnd}[1]{ \mathrm{Sym}({#1}) }
\newcommand{\AsymEnd}[1]{ \mathrm{Asym}({#1}) }

%
\newcommand{\GF}[1]{ \mathcal{G}_{#1} }
\newcommand{\GFd}[1]{ \mathcal{G}^{*}_{#1} }

% Sobolev spaces of varifolds
\newcommand{\VSobnorm}[3]{\mathbf{W}_{#1} ( {#2}, {#3} )}
\newcommand{\VSobloc}[1]{\mathbf{W}_{#1}^{\mathrm{loc}}}
\newcommand{\VSob}[1]{\mathbf{W}_{#1}}
\newcommand{\VSobz}[1]{\mathbf{W}^{\diamond}_{#1}}

% the L^q norm of a function u in U is '\norm{u}{q}{A}' for "|u|_{q;U}" with
% bold |
\newcommand{\norm}[3]{\boldsymbol{|} #1 \boldsymbol{|}_{#2;#3}}

% \Shuffle(p,q)
% Shuffle of type (p,q), i.e., the set of permutations of {1,2,...,p+q}
% which are increasing on each of the sets {1,...,p} and {p+1,...p+q};
% see 1.3.3 in Federer1969.
\newcommand{\Shuffle}[2]{\operatorname{\textbf{Sh}}(#1,#2)}

% \Choice(n,k)
% All choices of k elements from the set of n elements realised
% as the set of all increasing maps from {1,...,k} to {1,...,n};
% see 1.3.2 in Federer1969.
\newcommand{\Choice}[2]{\boldsymbol{\Lambda}(#1,#2)}

% \Partition(k,n) (a.k.a. multi-indices on \R^k of length n)
% Partitions of n into k subsets realised as the set of functions
% \alpha mapping the set {1,...,k} into non-negative integers
% such that \sum_{i=1}^k \alpha(i) = n ; see 1.9.2 in Federer1969.
\newcommand{\Partition}[2]{\boldsymbol{\Xi}(#1,#2)}

% \Sequence(n,k)
% The set of functions of the type {1,...,k} --> {1,...,n} ;
% see 1.10.6 in Federer1969.
\newcommand{\Sequence}[2]{\mathscr{S}(#1,#2)}

% The space of continuous functions with compact support, "K(X)" with
% calligraphic "K", produced by '\ccspace{X}'
\newcommand{\ccspace}[1]{\mathscr{K}(#1)}

% The space of smooth test functions with compact support
\newcommand{\dspace}[2]{\mathscr{D}(#1,#2)}

% The space of distributions
\newcommand{\distr}[2]{\mathscr{D}'(#1,#2)}

% The space of smooth functions
\newcommand{\espace}[2]{\mathscr{E}(#1,#2)}

% turning a plane into a projection, i.e. obtain '\project{T}' produces
% "T_\natural", the projection associated with the plane T
\newcommand{\project}[1]{{#1}_\natural}

% similar, with parentheses
\newcommand{\eqproject}[1]{(#1)_\natural}

% again similar for the orthogonal complement of the plane T
\newcommand{\perpproject}[1]{{#1}^\perp_\natural}

% the Grassmannian '\grass{n}{m}' produces "G(n,m)" with bold "G", i.e. the
% set of m dimensional planes in R^n
\newcommand{\grass}[2]{\mathbf{G}(#1,#2)}

% The projective Grassmannian,
% i.e., $\pgrass{X}{k}$ is the set of all orthogonal projections in X of rank k
\newcommand{\pgrass}[2]{\mathbf{G}_{\natural}(#1,#2)}

% the oriented Grassmannian
\newcommand{\ograss}[2]{\mathbf{G}_0(#1,#2)}

\newcommand{\AllP}{\mathbf{P}}

% the manifold of all projections
\newcommand{\Proj}[2]{\AllP(#1,#2)}

% The measure of a ball, '\measureball{\mu}{\oball{a}{r}}' produces '\mu
% U(a,r)', i.e. no brackets for measure of balls for readability
\newcommand{\measureball}[2]{{#1}\,{#2}}

% Space of Q valued functions
\newcommand{\qspace}{\mathbf{Q}}

% Commands for varifolds
\newcommand{\Var}[1]{\mathbf{V}_{#1}}     % general varifolds
\newcommand{\RVar}[1]{\mathbf{RV}_{#1}}   % rectifiable varifolds
\newcommand{\IVar}[1]{\mathbf{IV}_{#1}}   % integral varifolds
\DeclareMathOperator{\VarTan}{VarTan}     % varifold tangents
\newcommand{\var}[1]{\mathbf{v}_{#1}}     % varifolds generated by rectifiable sets

\DeclareMathOperator{\graph}{graph}     % the graph of a function

% For ease of changing these letters
\newcommand{\adim}{n}
\newcommand{\vdim}{k}
\newcommand{\codim}{n-k}

% integral currents
\newcommand{\IC}{\mathbf{I}}

% integral currents
\newcommand{\NC}{\mathbf{N}}

% the open ball: '\oball{a}{r}' produces "U(a,r)" with bold "U"
\newcommand{\oball}[2]{\mathbf{U}(#1,#2)}

% the closed ball: '\cball{a}{r}' produces "B(a,r)" with bold "B"
\newcommand{\cball}[2]{\mathbf{B}(#1,#2)}

% the closed ball: '\cball{a}{r}' produces "B(a,r)" with bold "B"
\newcommand{\cballa}[3]{\mathbf{B}^{#1}(#2,#3)}

% the sphere
\newcommand{\sphere}[1]{\mathbb{S}^{#1}}

% the canonical isomorphism between \End{X} and X^* \otimes X
\newcommand{\caniso}[1]{\boldsymbol{\Gamma}_{\!#1} }

% textstyle integral signs
\ifdefined\textint
    \renewcommand{\textint}[2]{{\textstyle\int_{#1}^{#2}}}
\else    
    \newcommand{\textint}[2]{{\textstyle\int_{#1}^{#2}}}
\fi

% textstyle integral signs
\ifdefined\textfint
    \renewcommand{\textfint}[2]{{\textstyle\fint_{#1}^{#2}}}
\else    
    \newcommand{\textfint}[2]{{\textstyle\fint_{#1}^{#2}}}
\fi

% textstyle sum sign
\ifdefined\textsum
    \renewcommand{\textsum}[2]{{\textstyle\sum_{#1}^{#2}}}
\else  
    \newcommand{\textsum}[2]{{\textstyle\sum_{#1}^{#2}}}
\fi

% textstyle product
\ifdefined\textint
    \renewcommand{\textprod}[2]{{\textstyle\prod_{#1}^{#2}}}
\else    
    \newcommand{\textprof}[2]{{\textstyle\prod_{#1}^{#2}}}
\fi


% natural numbers
% \newcommand{\nat}{\mathbb{N}}
% \newcommand{\nat}{\mathbb{N}_+}
\newcommand{\natp}{\mathscr{P}}
\newcommand{\nat}{\natp \cup \{0\}}

% the integers
\newcommand{\integers}{\mathbf{Z}}
\newcommand{\Z}{\integers}

% for 0-dimensional currents: [[ x ]]
% \newcommand{\Lcur}{\boldsymbol{[}}
% \newcommand{\Rcur}{\boldsymbol{]}}
% \newcommand{\Lcur}{[\![}
% \newcommand{\Rcur}{]\!]}
\newcommand{\Lcur}{\llbracket}
\newcommand{\Rcur}{\rrbracket}

% admissible deformations
\newcommand{\adm}[1]{\mathfrak{D}({#1})}

% complex numbers
\newcommand{\C}{\mathbf{C}}

% real numbers
\newcommand{\R}{\mathbf{R}}

% cubical subdivision of R^n
\newcommand{\W}{\mathbf{W}}

% characteristic function
% \newcommand{\CF}{\mathds{1}}
\newcommand{\CF}[1]{\boldsymbol{\chi}_{#1}}

% the Čech homology
\newcommand{\Ch}{\mathbf{\check H}}

% the singular homology
\newcommand{\Sh}{\mathbf{H}}

% homology
\newcommand{\h}{\mathbf{H}}

% cycles for homology theory
\newcommand{\hZ}{\mathscr{Z}}

% boundaries  for homology theory
\newcommand{\hB}{\mathscr{B}}

% the Lebesgue measure
\newcommand{\LM}{\mathscr{L}}

% The Lebesgue space, '\Lp{p}' produces "L_p" with bold L
\newcommand{\Lp}[1]{\mathbf{L}_{#1}}

% The local Lebesgue space, '\Lploc{p}' produces "L_p^{loc}" with bold L
\newcommand{\Lploc}[1]{\mathbf{L}_{#1}^{\mathrm{loc}}}

% Slightly old-fashened notation, may explain or replace
% The L^p-norm of u w.r.t. \phi, '\Lpnorm{\phi}{p}{u}' produces
% "\phi_{(p)}(u)", unusual but efficient. Also
\newcommand{\Lpnorm}[3]{{#1}_{({#2})}({#3})}
\newcommand{\eqLpnorm}[3]{{(#1)}_{({#2})}({#3})}

% the number in Besicovitch's covering theorem in n dimensions,
% '\besicovitch{n}', produces "\beta(n)" with bold beta
\newcommand{\besicovitch}[1]{\boldsymbol{\beta}(#1)}

% The nearest point projection
\newcommand{\npp}[1]{\boldsymbol{\xi}_{{#1}}}

% The distance function
% \newcommand{\dist}[1]{\boldsymbol{\delta}_{{#1}}}
\DeclareMathOperator{\dist}{dist}

% the distance function
\newcommand{\distF}[2]{\boldsymbol{\delta}_{#2}^{#1}}
% \da{A}{\phi} will produce \boldsymbol{\delta}_{A}^{\phi}
\newcommand{\da}[2]{\distF{#2}{#1}}

% nu
\newcommand{\an}[2]{\boldsymbol{\nu}_{#1}^{#2}}

% r
\newcommand{\ar}[2]{\boldsymbol{r}_{#1}^{#2}}

% \xi
\newcommand{\anp}[2]{\boldsymbol{\xi}_{#1}^{#2}}

% boldface p
\newcommand{\pp}{\mathbf{p}}

% boldface q
\newcommand{\qq}{\mathbf{q}}

% The domain of the unique nearest point projection
\newcommand{\Unp}[1]{\mathrm{Unp}({#1})}
% \DeclareMathOperator{\Unp}{Unp}


% % distance
% \DeclareMathOperator{\dist}{dist}

% the Hausdorff measure
\newcommand{\HM}{\mathscr{H}}

% the Hausdorff distance
\newcommand{\HD}{{d_{\HM}}}
% \newcommand{\HD}{\operatorname{HD}}

\newcommand{\HDK}[1]{{d_{\HM,#1}}}

% the flat metric on varifolds
\newcommand{\FM}{\mathbf{L}}

% the density function for measures
\newcommand{\density}{\boldsymbol{\Theta}}
\newcommand{\udensity}{\boldsymbol{\overline{\Theta}}}
\newcommand{\ldensity}{\boldsymbol{\underline{\Theta}}}

\newcommand{\BFconst}{\boldsymbol{\beta}}

\newcommand{\eqclsl}{\boldsymbol{[}}
\newcommand{\eqclsr}{\boldsymbol{]}}

% measure of the n dimensional unit ball, '\unitmeasure{n}' produces
% "\alpha(n)" with bold alpha
\newcommand{\unitmeasure}[1]{\boldsymbol{\alpha}(#1)}

% The restriction of a measure \phi to a set A is '\phi \restrict A'
% \newcommand{\restrict}{\mathop{\llcorner}}
\newcommand{\restrict}{ \mathop{ \rule[0pt]{.5pt}{6pt} \rule[0pt]{4pt}{0.5pt} }\nolimits }
% \newcommand{\restrict}{\lefthalfcup}

% The contraction of a form with a vectorfield
% \newcommand{\contract}{\righthalfcup}
\newcommand{\contract}{ \mathop{ \rule[0pt]{4pt}{0.5pt} \rule[0pt]{.5pt}{6pt} }\nolimits }

% the upright "d" in integrals
\newcommand{\ud}{\ensuremath{\,\mathrm{d}}}

% the upright "D" for differentiation
\newcommand{\uD}{\ensuremath{\mathrm{D}}}
% \DeclareMathOperator{\Der}{D}       
\newcommand{\Der}{\uD}

\DeclareRobustCommand{\rchi}{{\mathpalette\irchi\relax}}
\newcommand{\irchi}[2]{\raisebox{\depth}{$#1\chi$}}

\newcommand{\arl}[2]{\underline{\ar{#1}{#2}}}

\DeclareMathOperator{\cato}{cato}

\DeclareMathOperator{\extreme}{extreme}
\DeclareMathOperator{\exposed}{exposed}


% the weak derivative with bold D instead of D for the classical derivative
% \newcommand{\weakD}{\operatorname{\mathbf{D}}}
\newcommand{\weakD}{\ensuremath{\mathbf{D}}}

% the identity mapping on A, '\id{A}', produces "1_A" with bold 1
% \newcommand{\id}[1]{\mathbf{1}_{#1}}
% \newcommand{\id}[1]{\operatorname{id}_{#1}}
% \newcommand{\id}[1]{\mathrm{id}_{#1}}
% \newcommand{\id}[1]{\mathbf{I}_{#1}}
\newcommand{\id}[1]{\mathds{1}_{#1}}
% \newcommand{\id}[1]{\mathbb{1}_{#1}}

% The image of a set A under a map f is "f[A]"
\newcommand{\lIm}{[}
\newcommand{\rIm}{]}

\newcommand{\scale}[1]{\boldsymbol{\mu}_{#1}}
\newcommand{\trans}[1]{\boldsymbol{\tau}_{#1}}

% the space of a simplicial complex
\newcommand{\sspace}[1]{\textstyle{\bigcup}\,{#1}}

% \newcommand{\tbwedge}{\text{{\small $\bigwedge$}}}
\newcommand{\tbwedge}{{\textstyle\boldsymbol{\bigwedge}}}

% exterior algebra
% \newcommand{\extpower}[1]{\text{\raisebox{2pt}{\scriptsize $\bigwedge$}}_{#1}}
% \newcommand{\extpower}[1]{{\textstyle \bigwedge^{#1}}}

% space of anti-symmetric forms
\newcommand{\aforms}[1]{\tbwedge^{\!{#1}}}
% \newcommand{\aforms}[1]{\bigwedge\nolimits^{\!{#1}}}

% exterior power of a vectorspace
\newcommand{\extpower}[1]{\tbwedge_{{#1}}}
% \newcommand{\extpower}[1]{\bigwedge\nolimits_{{#1}}}

% % exterior algebra
% % \newcommand{\extpower}[1]{\text{\raisebox{2pt}{\scriptsize $\bigwedge$}}_{#1}}
% \newcommand{\extpower}[1]{{\textstyle \bigwedge^{#1}}}
% \newcommand{\tbwedge}{{\textstyle \bigwedge}}

% Adjoint of a linear map
\newcommand{\adjoint}[1]{\aforms{1} {#1}}
\newcommand{\aadjoint}[1]{\aforms{1} \aforms{1} {#1}}

% transpose of a linear map between Euclidean spaces
% \transpose{A} produces A^{*}
\newcommand{\transpose}[1]{ #1^{*} }
\newcommand{\ttranspose}[1]{ #1^{**} }

% convex transpose of a convex map
\newcommand{\ctran}[1]{{#1}^{\dagger}}
\newcommand{\cttran}[1]{{#1}^{\dagger\dagger}}

% polar of a function or a set
\newcommand{\polar}[1]{{#1}^{\circ}}
\newcommand{\ppolar}[1]{{#1}^{\circ\circ}}

% convex conjugate of a extended real valued map
\newcommand{\cconj}[1]{{#1}^{\star}}
\newcommand{\ccconj}[1]{{#1}^{\star\star}}

% textstyle big cup
\newcommand{\tbcup}{{{\textstyle \bigcup}}}

% textstyle big cap
\newcommand{\tbcap}{{{\textstyle \bigcap}}}

% The closure of a set A, '\Clos{A}' produces "Clos A"
\newcommand{\Clos}[1]{\mathop{\mathrm{Clos}}#1}

% The centre, side length, and "lower left" corner of a cube
\newcommand{\side}[1]{\mathbf{l}(#1)}
\newcommand{\centre}[1]{\mathbf{c}(#1)}
\newcommand{\vertex}[1]{\mathbf{o}(#1)}

% The interior and boundary of a cube
\newcommand{\cInt}[1]{\operatorname{Int}_{\mathrm{c}}(#1)}
% \newcommand{\cBdry}[1]{\operatorname{Bdry}_{\mathrm{c}}(#1)}
\newcommand{\cBdry}[1]{\partial_{\mathrm{c}}#1}
% \newcommand{\cInt}[1]{\mathop{\mathbf i}(#1)}
% \newcommand{\cBdry}[1]{\mathop{\mathbf b}(#1)}

% The standard dyadic cubes
\newcommand{\cubes}{\mathbf{K}}

% The Whitney family of cubes
\newcommand{\WF}{\mathbf{WF}}

% The cubical complex of a family of cubes
\newcommand{\CX}{\mathbf{CX}}

% Two commands for cylinders, '\cylind{T}{a}{r}' produces "C(T,a,r)" with bold
% C, for plane T, center a, radius r
\newcommand{\cylind}[3]{\mathbf{C} ( #1, #2, #3 )}

% a variant for cylinders with finite height
\newcommand{\cylinder}[4]{\mathbf{C} ( #1, #2, #3, #4 )}

% A symbol for the Dirac measure
% \newcommand{\Dirac}[1]{\operatorname{Dirac}(#1)}
% ... following the French/Bourbaki school
% \newcommand{\Dirac}[1]{\varepsilon}_{#1}
% ... following Federer's 1969 book 2.5.19
\newcommand{\Dirac}[1]{ \boldsymbol{\delta}_{#1} }

% Smooth compactly supported vectorfields
\newcommand{\VF}{\mathscr{X}}

% A symbol for the Dirac measure
\newcommand{\ellipticity}[1]{\operatorname{ellipticity}(#1)}

% affine approximate
% \DeclareMathOperator{\Aff}{A}       

% affine space generated by a set
\DeclareMathOperator{\aff}{af{}f}

% The trace
\DeclareMathOperator{\trace}{trace}

\DeclareMathOperator{\Tr}{Tr}

% The greatest commands divisor
% \DeclareMathOperator{\gcd}{gcd}

% The neighbours of a cube
% \DeclareMathOperator{\nbs}{nbs}

% Space of homomorphisms between vector spaces
\DeclareMathOperator{\Hom}{Hom}

% The boundary of a set
\DeclareMathOperator{\Bdry}{Bdry}

% something approximate
\DeclareMathOperator{\ap}{ap}

\DeclareMathOperator*{\aplim}{\mathrm{ap}\, \lim}   % approximate limit
\DeclareMathOperator*{\aplimsup}{\mathrm{ap}\, \limsup}   % approximate limit

% linear span
\DeclareMathOperator{\lin}{span}

% class of continuity
\newcommand{\cnt}[1]{\mathscr{C}^{#1}}

% skeleton of a simplicial complex
\newcommand{\skel}[2]{\mathrm{sk}_{#1}(#2)}

% the set of orthogonal projections "O^*(n,m)" from R^n onto R^m with bold "O",
% produced by '\orthproj{n}{m}'
\newcommand{\orthproj}[2]{\mathbf{O}^\ast({#1},{#2})}
\newcommand{\orthinj}[2]{\mathbf{O}({#1},{#2})}

\newcommand{\powerset}[1]{\mathbf{2}^{#1}}

% the orthogonal group O(n)
\newcommand{\orthgroup}[1]{\mathbf{O}({#1})}

% the H\"older semi norm of f with exponent \alpha is "h_{\alpha} (f)" with
% bold \alpha, '\hoelder{\alpha}{f}'
\newcommand{\hoelder}[2]{\mathbf{h}_{#1}(#2)}

% the evaluation map  \ev{X} : X --> X^{**}
\newcommand{\ev}[1]{\mathbf{ev}_{#1}}

% An unnamed parameter/placeholder
\newcommand{\ph}{\,\cdot\,}

% Interior of a set
\DeclareMathOperator{\Int}{Int}

% Relative interior of a set
\DeclareMathOperator{\rInt}{relint}

% Barrier cone
\DeclareMathOperator{\bc}{bc}

% Effective domain of a convex function
\DeclareMathOperator{\dom}{dom}

% Relative boundary of a set
\DeclareMathOperator{\rBdry}{relbd}

% domain
\DeclareMathOperator{\reach}{reach}

% domain
\DeclareMathOperator{\dmn}{dmn}

% gradient
\DeclareMathOperator{\grad}{grad}
% \DeclareMathOperator{\grad}{\nabla}

% subgradient of a convex function
% \newcommand{\subgrad}{\text{{\large $\nabla$}}}
% \newcommand{\subgrad}{\nabla}
\DeclareMathOperator{\subgrad}{subgr}

% support
\DeclareMathOperator{\spt}{spt}

% space generated be a k-vector
\DeclareMathOperator{\Sp}{space}

% convex hull
\DeclareMathOperator{\conv}{conv}

% uniform convex hull
\DeclareMathOperator{\uconv}{uc}

% cone over a set (as in the book of Rockafellar)
\DeclareMathOperator{\ray}{ray}

% convex cone over a set (as in the book of Rockafellar)
\DeclareMathOperator{\cone}{cone}

% tangent space
\DeclareMathOperator{\Tan}{Tan}

% normal space
\DeclareMathOperator{\Nor}{Nor}

% Lipschitz constant
\DeclareMathOperator{\Lip}{Lip}

% produces "card", cardinality of set
\DeclareMathOperator{\card}{card}

% diameter
\DeclareMathOperator{\diam}{diam}

% sign of
\DeclareMathOperator{\sgn}{sgn}

% % maximum
% \DeclareMathOperator{\max}{max}

% Set A without B, 'A \without B' produces "A \sim B"
\newcommand{\without}{\!\sim\!}
% \DeclareMathOperator{\without}{\sim}

\newcommand{\QEC}{\textrm{QEC}}
\newcommand{\EC}{\textrm{EC}}
\newcommand{\SAC}{\textrm{SAC}}
\newcommand{\SEC}{\textrm{SEC}}
\newcommand{\CI}{\textrm{CI}}
\newcommand{\USAC}{\textrm{USAC}}
\newcommand{\AC}{\textrm{AC}}
\newcommand{\AlmE}{\textrm{AE}}
\newcommand{\AlmUE}{\textrm{UAE}}

% classification { x : properties }, : is '\with'
\DeclareMathOperator{\with}{:}

% Subspace of End(R^n) with image in a fixed space
\newcommand{\imsp}{\mathcal{I}}

% image of a map
\DeclareMathOperator{\im}{im}

% boundary of a simplex
\DeclareMathOperator{\sBdry}{b}

% interior of a simplex
\DeclareMathOperator{\sInt}{i}

% barycentre of a simplex
\DeclareMathOperator{\baryc}{c}

% AUE
\DeclareMathOperator{\FAUE}{AUE}
% AE
\DeclareMathOperator{\FAE}{AE}
% BC
\DeclareMathOperator{\FBC}{BC}
% wBC
\DeclareMathOperator{\FwBC}{wBC}
% AC
\DeclareMathOperator{\FAC}{AC}

% To mark parapraph which lack final verification
\newcommand{\verify}{\framebox{{\color{Red} \textsc{verify}}~}}
\newcommand{\todo}{\framebox{{\color{Red} \textsc{to do}}}}
\newcommand{\maciej}[1]{\begin{center}\framebox{\begin{minipage}{0.95\textwidth} {\color{Fuchsia}\textsc{@ML:} \textsf{#1}} \end{minipage}}\end{center}}
\newcommand{\comment}[1]{{\color{teal}$\boldsymbol{\gg}$ \textsf{#1} $\boldsymbol{\ll}$}}

\newcommand{\cred}{\color{red}}
\newcommand{\cg}{\color{green}}
\newcommand{\cm}{\color{magenta}}


% Local Variables:
% coding: utf-8
% eval: (custom-set-variables  '(LaTeX-indent-level 4) '(LaTeX-item-indent -4) '(indent-tabs-mode nil))
% eval: (ispell-change-dictionary "british")
% mode: latex
% eval: (flyspell-mode)
% End:


% Local notation of this note.
\newcommand{\heatop}{\mathcal{H}}          % heat operator  \partial_t + \Delta
\newcommand{\sv}[1]{\varsigma_{#1}}        % the numbers of 2.4
\newcommand{\dUps}{\overrightarrow{\Upsilon}} % the polar of Upsilon, Federer 4.1.5
\setlength{\emergencystretch}{2em}

\date{\today}

\title{Blueprint of a formalisation of the anisotropic Michael--Simon inequality of De~Philippis,
    Gennaioli, Pigati and Rindler}

\author{Working note}

\begin{document}

\maketitle

\tableofcontents

\section{Introduction}
\label{sec:intro}

De~Philippis, Gennaioli, Pigati and Rindler~\cite[Theorem~1.1]{DePhilippisGennaioliPigatiRindler2026}
prove a scale-invariant bound on the mass of a~measure satisfying a~first order constraint. Suppose
$\mu = (\HM^{\vdim} \restrict M) \restrict \theta$ is a~finite measure carried by a~Borel set $M
\subseteq \R^{\adim}$ with $\HM^{\vdim}(M) < \infty$. Suppose a~measurable map $X$ takes values in
a~compact convex set $K$ of linear maps whose $\vdim$-th singular values are bounded away from
zero, and the divergence of the matrix valued measure $X\mu$ is a~finite measure. Then
$\mu(\R^{\adim}) \le C \, \HM^{\vdim}(M)^{1/\vdim} \, |\mathrm{div}(X\mu)|(\R^{\adim})$. As
an~application, the same authors obtain the Michael--Simon inequality for rectifiable varifolds
whose first variation with respect to an~anisotropic integrand~$F$ is a~finite measure, provided
$F$ satisfies the condition (AC1) of De~Philippis, De~Rosa and
Ghiraldin~\cite[Theorem~1.6]{DePhilippisGennaioliPigatiRindler2026}. The proof applies the heat
semigroup to a~part of~$\mu$ carried by a~compact rectifiable set. It bounds the Fisher
information of the heated measure by the constraint, compares the resulting growth of its entropy
with the Minkowski content of the compact set, and adapts the heat-flow identities of Gennaioli
and Rindler~\cite[section~3]{Gennaioli2026}.

This note is a~blueprint for a~formalisation of these two results, in Lean or in Rocq. Each
assertion of the proof is matched with an~item of Federer's book~\cite{Federer1969}, an~item of
Allard's paper~\cite{Allard1972}, a~lemma of De~Philippis, De~Rosa and
Ghiraldin~\cite{DePhilippis2018}, or a~lemma of this note. Every lemma of this note is stated with
all its hypotheses and proved from the preceding items. The notation is that of Federer.
Vector valued measures are replaced by linear functionals on spaces of continuous maps, as
in~\cite[2.5.12, 2.5.14, 4.1.5]{Federer1969}, and varifolds are treated as
in~\cite{Allard1972}. Hence the lemmas of this note, together with~\cite[Lemmas~A.1
and~A.2]{DePhilippis2018}, are the results to be formalised on top of formalisations of
Federer's book and of Allard's paper.

Theorem~\ref{thm:main} reproduces~\cite[Theorem~1.1]{DePhilippisGennaioliPigatiRindler2026} under
the additional hypothesis that $M$ is countably $(\HM^{\vdim},\vdim)$ rectifiable. The authors
of~\cite{DePhilippisGennaioliPigatiRindler2026} remove this hypothesis
by~\cite[Proposition~3.1]{ArroyoRabasa2019}; see~\cite[Remark~1.2 and the proof of
Theorem~1.1]{DePhilippisGennaioliPigatiRindler2026}. The additional hypothesis holds in the
application to varifolds. Hence Theorem~\ref{thm:MS}
reproduces~\cite[Theorem~1.6]{DePhilippisGennaioliPigatiRindler2026} in full. Theorems~1.3
and~1.4 of~\cite{DePhilippisGennaioliPigatiRindler2026} are not treated.

The deduction follows~\cite{DePhilippisGennaioliPigatiRindler2026} line by line, with four
changes. First, the cut-off arguments mentioned
in~\cite{DePhilippisGennaioliPigatiRindler2026} for integrations by parts on~$\R^{\adim}$ are
replaced by the classes of functions of polynomial growth and of Gaussian decay
of~\ref{def:classes}. Second, the smooth cut-off function~$\chi$ of~\cite[proof of
Lemma~3.2]{DePhilippisGennaioliPigatiRindler2026} is replaced by a~piecewise affine one. This
requires only the Lipschitz property~\ref{lem:spectral} of the spectral calculus, instead of its
smoothness, and the integration by parts~\ref{lem:ibp} for locally Lipschitzian vector fields.
Third, the singular values are replaced by the numbers $\sv{\vdim}(X)$ of~\ref{def:sv}, and
$\Lambda$ bounds the Hilbert--Schmidt norm, as assumed at the beginning of~\cite[section~3]
{DePhilippisGennaioliPigatiRindler2026}. Fourth, the Minkowski content estimate in the proof
of~\cite[(4.2)]{DePhilippisGennaioliPigatiRindler2026} is deduced from~\cite[3.2.39]{Federer1969}
for an~arbitrary compact $\vdim$ rectifiable set. Table~\ref{tab:correspondence} lists the
correspondence of the items. Table~\ref{tab:dependencies} in section~\ref{sec:summary} lists the
dependencies of all items of this note.

\begin{table}[h]
    \centering
    \begin{tabular}{p{0.52\textwidth}p{0.3\textwidth}}
        \cite{DePhilippisGennaioliPigatiRindler2026} & this note \\ \hline
        (3.1), (3.2), (3.4) & \ref{def:heat_objects} \\
        (3.3), (3.4), the identity before (3.12) & \ref{lem:identities} \\
        (3.6) and the identities in the proof of Lemma~3.1 & \ref{lem:heat_calculus} \\
        Lemma~3.1 with (3.7)--(3.12) & \ref{lem:DGPR31} \\
        (3.14), (3.15), (3.18) & proof of~\ref{lem:DGPR32} \\
        (3.16) & \ref{lem:semiconvex} \\
        (3.17) & \ref{lem:ibp} \\
        Lemma~3.2 & \ref{lem:DGPR32} \\
        (4.1) & \ref{lem:upper_entropy} \\
        (4.2) & \ref{lem:tube}, \ref{lem:lower_entropy} \\
        proof of Theorem~1.1 & \ref{lem:inner}, \ref{lem:compact_case}, \ref{thm:main} \\
        first variation formula of section~1.2 & \ref{lem:first_variation} \\
        Theorem~1.6 & \ref{thm:MS}
    \end{tabular}
    \caption{Items of~\cite{DePhilippisGennaioliPigatiRindler2026} and of this note.}
    \label{tab:correspondence}
\end{table}

\section{Notation and auxiliary linear algebra}
\label{sec:notation}

\begin{miniremark}
    \label{mr:sec_notation}
    This section fixes the notation and recalls linear functionals representable by integration.
    It then introduces the numbers $\sv{\vdim}(X)$ and proves five auxiliary facts: properties of
    $\sv{\vdim}$ in~\ref{lem:sv}, a~Lipschitzian spectral calculus in~\ref{lem:spectral},
    a~trace comparison in~\ref{lem:trace}, compactness of convex hulls in~\ref{lem:hull}, and the
    position of barycentres in~\ref{lem:barycentre}.
\end{miniremark}

\begin{miniremark}
    \label{mr:notation}
    Throughout, $\adim$ and $\vdim$ are integers with $1 \le \vdim \le \adim$, and $Y$ is
    a~finite dimensional inner product space. We fix an~orthonormal basis $\varepsilon_1, \ldots,
    \varepsilon_{\adim}$ of~$\R^{\adim}$ and its dual basis $\omega_1, \ldots, \omega_{\adim}$
    of~$\Hom(\R^{\adim},\R)$. The quantities below defined through these bases are traces or norms
    in the sense of~\cite[1.7.9, 1.7.10]{Federer1969}; hence they do not depend on the choice of
    the bases.

    Suppose $V$ and $Z$ are finite dimensional inner product spaces. Following~\cite[1.7.4,
    1.7.9]{Federer1969}, $\transpose{A} \in \Hom(Z,V)$ is the adjoint of $A \in \Hom(V,Z)$, the
    space $\Hom(V,Z)$ carries the inner product $A \bullet B = \trace(\transpose{A} \circ B)$, and
    $|A| = (A \bullet A)^{1/2}$. We write $\|A\| = \sup \{ |A(v)| : v \in V,\ |v| \le 1 \}$. Then
    $\|A\| \le |A| \le (\dim V)^{1/2} \|A\|$ by~\cite[1.7.9]{Federer1969}, $\|\transpose{A}\| =
    \|A\|$ by~\cite[1.7.6]{Federer1969}, and $|B \circ A| \le \|B\| \, |A|$ and $|A \circ C| \le |A| \, \|C\|$ whenever $B \in
    \Hom(Z,Z')$ and $C \in \Hom(V',V)$. For $\psi \in \Hom(Z,\R)$ and $z \in Z$ we write $\langle z, \psi \rangle =
    \psi(z)$ and $\|\psi\| = \sup \{ \langle z, \psi \rangle : |z| \le 1 \}$. For $\omega \in
    \Hom(V,\R)$ and $b \in Z$, the map $\omega \cdot b \in \Hom(V,Z)$ is given by $(\omega \cdot
    b)(v) = \omega(v) b$ for $v \in V$; then $(\omega \cdot b) \bullet A = \omega(\transpose{A}(b))$
    for $A \in \Hom(V,Z)$; see~\cite[2.3(4)]{Allard1972}. We write $\SymEnd{\R^{\adim}} =
    \End{\R^{\adim}} \cap \{ A : \transpose{A} = A \}$. Each $A \in \SymEnd{\R^{\adim}}$ admits an
    orthonormal basis of~$\R^{\adim}$ consisting of eigenvectors of~$A$
    by~\cite[1.7.3]{Federer1969}.

    Suppose $Z$ is a~finite dimensional normed space. Following~\cite[1.10.1,
    3.1.11]{Federer1969}, $\textstyle\bigodot^j(\R^{\adim},Z)$ is the space of symmetric $j$~linear
    maps of $\R^{\adim}$ into~$Z$, and $\uD^j f(x) \in \textstyle\bigodot^j(\R^{\adim},Z)$ is the
    $j$-th differential at~$x$ of a~map~$f$ of class~$j$. For $f : \R^{\adim} \to \R$ of class~$2$
    we define $\grad f(x) \in \R^{\adim}$ by $\uD f(x)(v) = \grad f(x) \bullet v$ for $v \in
    \R^{\adim}$, and
    \begin{displaymath}
        \Delta f(x) = \textsum{j=1}{\adim} \uD^2 f(x)(\varepsilon_j,\varepsilon_j) \,;
    \end{displaymath}
    the same formula defines $\Delta f(x) \in Z$ for $f : \R^{\adim} \to Z$ of class~$2$. For $g :
    \R^{\adim} \to \R^{\adim}$ and $W : \R^{\adim} \to \Hom(\R^{\adim},Y)$ of class~$1$ we define
    $\mathrm{div}\, g(x) = \trace \uD g(x) = \textsum{j=1}{\adim} \uD g(x)(\varepsilon_j) \bullet
    \varepsilon_j$, the covector $\mathrm{div}\, W(x) \in \Hom(Y,\R)$ by
    \begin{displaymath}
        \langle y, \mathrm{div}\, W(x) \rangle = \textsum{j=1}{\adim} \uD W(x)(\varepsilon_j)
        (\varepsilon_j) \bullet y
        = \trace \bigl[ \R^{\adim} \ni v \mapsto \transpose{(\uD W(x)(v))}(y) \bigr]
    \end{displaymath}
    for $y \in Y$,
    and $|\uD W(x)|^2 = \textsum{j=1}{\adim} |\uD W(x)(\varepsilon_j)|^2$, which is the norm
    of~\cite[1.7.9]{Federer1969} on~$\Hom(\R^{\adim},\Hom(\R^{\adim},Y))$.

    For $0 < T < \infty$ we set $\Omega_T = (\R \times \R^{\adim}) \cap \{ (t,x) : t < T \}$. For
    a~map $f$ defined on~$\Omega_T$ we write $f_t = f(t,\cdot)$; then $\uD f_t(x)$, $\uD^2 f_t(x)$,
    $\grad f_t(x)$, $\Delta f_t(x)$ are taken with respect to~$x$, $\partial_t f(t,x)$ is the
    derivative with respect to~$t$, and $\uD^j f(t,x)$ is the $j$-th differential with respect to
    $(t,x) \in \R \times \R^{\adim}$.

    We use $\LM^{\adim}$, $\HM^{\vdim}$, $\cball{a}{r}$, $\oball{a}{r}$, $\dist(x,E)$,
    $\unitmeasure{m} = \LM^m(\cball{0}{1})$, $\phi \restrict A$, and the measure $\phi \restrict
    f$ with density~$f$ as in~\cite[2.4.10]{Federer1969}. By~\cite[2.7.16(1)]{Federer1969},
    $\unitmeasure{m} = \pi^{m/2} / \Gamma(1 + m/2)$. Following~\cite[2.5.14, 4.1.1]{Federer1969},
    $\ccspace{\R^{\adim}}$, $\ccspace{\R^{\adim}}^+$, $\ccspace{\R^{\adim},Z}$,
    $\dspace{\R^{\adim}}{Z}$ and $\distr{\R^{\adim}}{Z}$ denote the spaces of continuous functions
    with compact support, their nonnegative members, continuous maps into~$Z$ with compact
    support, smooth maps into~$Z$ with compact support, and distributions of type~$Z$. Following
    \cite[2.5]{Allard1972}, $\VF(\R^{\adim}) = \dspace{\R^{\adim}}{\R^{\adim}}$. The notions
    \emph{$\vdim$~rectifiable} and \emph{countably $(\HM^{\vdim},\vdim)$ rectifiable} are those
    of~\cite[3.2.14]{Federer1969}. We call a~measure $\phi$ over~$\R^{\adim}$ a~\emph{Borel
    measure} if all Borel subsets of~$\R^{\adim}$ are $\phi$~measurable.
\end{miniremark}

\begin{miniremark}
    \label{mr:functionals}
    Suppose $Z$ is a~finite dimensional inner product space and $\Theta : \ccspace{\R^{\adim},Z}
    \to \R$ is a~linear functional such that
    \begin{displaymath}
        \|\Theta\|(f) = \sup \bigl\{ \Theta(\xi) : \xi \in \ccspace{\R^{\adim},Z},\ |\xi(x)| \le
        f(x) \text{ for } x \in \R^{\adim} \bigr\} < \infty
        \quad \text{for $f \in \ccspace{\R^{\adim}}^+$} \,.
    \end{displaymath}
    By~\cite[2.5.12]{Federer1969} in the form for~$\ccspace{\R^{\adim}}$ given in the last
    paragraph of~\cite[2.5.14]{Federer1969}, $\|\Theta\|$ is a~monotone Daniell integral; we
    denote by~$\|\Theta\|$ also the associated Radon measure over~$\R^{\adim}$ and call it the
    variation measure of~$\Theta$. There is a~$\|\Theta\|$ measurable map $\overrightarrow{\Theta}
    : \R^{\adim} \to \Hom(Z,\R)$ such that $\|\overrightarrow{\Theta}(x)\| = 1$ for
    $\|\Theta\|$~almost all~$x$ and
    \begin{displaymath}
        \Theta(\xi) = \textint{}{} \langle \xi(x), \overrightarrow{\Theta}(x) \rangle
        \ud \|\Theta\| x
        \quad \text{for $\xi \in \ccspace{\R^{\adim},Z}$} \,.
    \end{displaymath}
    A~distribution $\Theta \in \distr{\R^{\adim}}{Z}$ is \emph{representable by integration} if
    $\|\Theta\|(f) < \infty$ for $f \in \ccspace{\R^{\adim}}^+$, where $\|\Theta\|(f)$ is defined
    with $\dspace{\R^{\adim}}{Z}$ in place of~$\ccspace{\R^{\adim},Z}$; the preceding statements
    then hold with $\dspace{\R^{\adim}}{Z}$ in place of~$\ccspace{\R^{\adim},Z}$
    by~\cite[4.1.5]{Federer1969}. If moreover $\|\Theta\|(\R^{\adim}) < \infty$, we put
    $\Theta(g) = \textint{}{} \langle g(x), \overrightarrow{\Theta}(x) \rangle \ud \|\Theta\| x$
    for every bounded continuous map $g : \R^{\adim} \to Z$.

    Two consequences of these definitions are used. If $\nu$ is a~Radon measure over~$\R^{\adim}$
    and $|\Theta(g)| \le \textint{}{} |g| \ud \nu$ for $g \in \dspace{\R^{\adim}}{Z}$, then
    $\Theta$ is representable by integration and $\|\Theta\|(A) \le \nu(A)$ for $A \subseteq
    \R^{\adim}$, because $\|\Theta\|(f) \le \textint{}{} f \ud \nu$ for $f \in
    \ccspace{\R^{\adim}}^+$ and both measures are Radon measures; see~\cite[2.5.14]{Federer1969}.
    If $C \subseteq \R^{\adim}$ is closed and $\Theta(g) = 0$ whenever $g \in
    \dspace{\R^{\adim}}{Z}$ and $\spt g \cap C = \varnothing$, then $\|\Theta\|(f) = 0$ for $f \in
    \ccspace{\R^{\adim}}^+$ with $\spt f \cap C = \varnothing$; hence $\|\Theta\|(\R^{\adim}
    \without C) = 0$ by~\cite[2.5.14]{Federer1969}.
\end{miniremark}

\begin{definition}
    \label{def:sv}
    For $X \in \Hom(\R^{\adim},Y)$ we define
    \begin{displaymath}
        \sv{\vdim}(X) = \sup \bigl\{ \inf \{ |X(v)| : v \in P,\ |v| = 1 \} : P \in
        \grass{\adim}{\vdim} \bigr\} \,,
    \end{displaymath}
    where $\grass{\adim}{\vdim}$ is the set of $\vdim$~dimensional linear subspaces
    of~$\R^{\adim}$.
\end{definition}

\begin{lemma}
    \label{lem:sv}
    Suppose $X, X' \in \Hom(\R^{\adim},Y)$ and $c \in \R$. Then the following four statements
    hold.
    \begin{enumerate}
        \item \label{i:sv:hom} $\sv{\vdim}(cX) = |c| \, \sv{\vdim}(X)$.
        \item \label{i:sv:lip} $|\sv{\vdim}(X) - \sv{\vdim}(X')| \le \|X - X'\|$.
        \item \label{i:sv:rank} $\sv{\vdim}(X) > 0$ if and only if $\dim \im X \ge \vdim$.
        \item \label{i:sv:eigen} If $0 < a \le \sv{\vdim}(X)$, $u_1, \ldots, u_{\adim}$ is an
        orthonormal basis of~$\R^{\adim}$, and $\transpose{X} \circ X (u_i) = a_i u_i$ for $i = 1,
        \ldots, \adim$, then $\card \{ i : a_i \ge a^2 \} \ge \vdim$.
    \end{enumerate}
\end{lemma}

\begin{proof}
    Statement~\ref{i:sv:hom} is evident. Since $|X(v)| \le |X'(v)| + \|X - X'\|$ for $|v| = 1$, we
    obtain $\sv{\vdim}(X) \le \sv{\vdim}(X') + \|X - X'\|$; exchanging $X$ and~$X'$
    gives~\ref{i:sv:lip}.

    Suppose $\dim \im X \ge \vdim$. Since $\dim (\ker X)^\perp = \dim \im X$, there is $P \in
    \grass{\adim}{\vdim}$ with $P \subseteq (\ker X)^\perp$. The function $[v \mapsto |X(v)|]$ is
    positive and continuous on the compact set $P \cap \{ v : |v| = 1 \}$; hence its infimum there
    is positive and $\sv{\vdim}(X) > 0$. Suppose $\dim \im X < \vdim$. Then $\dim \ker X > \adim -
    \vdim$, hence every $P \in \grass{\adim}{\vdim}$ contains a~unit vector of~$\ker X$, and
    $\sv{\vdim}(X) = 0$. This proves~\ref{i:sv:rank}.

    Suppose $0 < a' < a$. There is $P \in \grass{\adim}{\vdim}$ with $|X(v)| \ge a' |v|$ for $v \in
    P$. Let $L$ be the linear span of $\{ u_i : a_i < a'^2 \}$. For $v \in L$ we have $|X(v)|^2 =
    v \bullet \transpose{X} \circ X(v) = \textsum{i}{} a_i (u_i \bullet v)^2$, where $i$ ranges over
    $\{ i : a_i < a'^2 \}$; hence $|X(v)| < a' |v|$ if $v \ne 0$. Thus $P \cap L = \{0\}$, $\dim L
    \le \adim - \vdim$, and $\card \{ i : a_i \ge a'^2 \} \ge \vdim$. Since the set $\{ a_1,
    \ldots, a_{\adim} \}$ is finite, letting $a'$ increase to~$a$ yields~\ref{i:sv:eigen}.
\end{proof}

\begin{remark}
    \label{rem:sv_singular}
    Suppose $X \in \Hom(\R^{\adim},Y)$, $u_1, \ldots, u_{\adim}$ is an orthonormal basis
    of~$\R^{\adim}$, $\transpose{X} \circ X(u_i) = a_i u_i$ for $i = 1, \ldots, \adim$, and $a_1
    \ge a_2 \ge \cdots \ge a_{\adim}$; such a~basis exists by~\cite[1.7.3]{Federer1969}. Then
    $\sv{\vdim}(X) = a_{\vdim}^{1/2}$. Indeed, $|X(v)|^2 \ge a_{\vdim} |v|^2$ for $v$ in the linear
    span of $u_1, \ldots, u_{\vdim}$, whence $\sv{\vdim}(X) \ge a_{\vdim}^{1/2}$; and
    $\sv{\vdim}(X)^2 > a_{\vdim}$ would give $\card \{ i : a_i > a_{\vdim} \} \ge \vdim$
    by~\ref{lem:sv}\ref{i:sv:eigen} applied with some $a$ satisfying $a_{\vdim} < a^2 <
    \sv{\vdim}(X)^2$, which is impossible. Hence $\sv{\vdim}(X)$ is the $\vdim$-th singular value
    $\sigma_{\vdim}(X)$ of~\cite[Theorem~1.1]{DePhilippisGennaioliPigatiRindler2026}. This
    identification is used only to compare the constants of this note with those
    of~\cite{DePhilippisGennaioliPigatiRindler2026}.
\end{remark}

\begin{definition}
    \label{def:spectral}
    Suppose $\chi : \R \to \R$ is Lipschitzian, $A \in \SymEnd{\R^{\adim}}$, $u_1, \ldots,
    u_{\adim}$ is an orthonormal basis of~$\R^{\adim}$, and $A(u_i) = a_i u_i$ for $i = 1, \ldots,
    \adim$. We define $\hat\chi(A) \in \End{\R^{\adim}}$ by
    \begin{displaymath}
        \hat\chi(A)(v) = \textsum{i=1}{\adim} \chi(a_i) (u_i \bullet v) u_i
        \quad \text{for $v \in \R^{\adim}$} \,.
    \end{displaymath}
    The map $\hat\chi(A)$ does not depend on the choice of $u_1, \ldots, u_{\adim}$: for each $a
    \in \R$, the map $[v \mapsto \textsum{i}{} (u_i \bullet v) u_i]$, where $i$ ranges over $\{ i :
    a_i = a \}$, is the orthogonal projection of~$\R^{\adim}$ onto $\ker (A - a \id{\R^{\adim}})$.
\end{definition}

\begin{lemma}
    \label{lem:spectral}
    Suppose $\chi$ and $\hat\chi$ are as in~\ref{def:spectral}, and $A, A' \in
    \SymEnd{\R^{\adim}}$. Then $\hat\chi(A) \in \SymEnd{\R^{\adim}}$, every eigenvalue
    of~$\hat\chi(A)$ belongs to $\chi[\{ a : \ker(A - a \id{\R^{\adim}}) \ne \{0\} \}]$, and
    \begin{displaymath}
        |\hat\chi(A) - \hat\chi(A')| \le \Lip(\chi) \, |A - A'| \,.
    \end{displaymath}
\end{lemma}

\begin{proof}
    The first two assertions are evident from~\ref{def:spectral}. Let $u_1, \ldots, u_{\adim}$ and
    $a_1, \ldots, a_{\adim}$ be as in~\ref{def:spectral}, and $u'_1, \ldots, u'_{\adim}$ and
    $a'_1, \ldots, a'_{\adim}$ the corresponding objects for~$A'$. Let $\pi_i$ and $\pi'_j$ be the
    orthogonal projections of~$\R^{\adim}$ onto the linear spans of $\{u_i\}$ and $\{u'_j\}$. Then
    $\textsum{i}{} \pi_i = \id{\R^{\adim}} = \textsum{j}{} \pi'_j$, $A = \textsum{i}{} a_i \pi_i$,
    and $A' = \textsum{j}{} a'_j \pi'_j$. Hence
    \begin{displaymath}
        \hat\chi(A) - \hat\chi(A') = \textsum{i,j}{} (\chi(a_i) - \chi(a'_j)) \, \pi_i \circ \pi'_j
        \,, \qquad
        A - A' = \textsum{i,j}{} (a_i - a'_j) \, \pi_i \circ \pi'_j \,.
    \end{displaymath}
    Since $\pi_i \circ \pi_{i'} = 0$ for $i \ne i'$, $\pi'_j \circ \pi'_{j'} = 0$ for $j \ne j'$,
    and the trace is invariant under cyclic permutations,
    \begin{displaymath}
        (\pi_i \circ \pi'_j) \bullet (\pi_{i'} \circ \pi'_{j'})
        = \trace(\pi'_j \circ \pi_i \circ \pi_{i'} \circ \pi'_{j'}) = 0
        \quad \text{whenever $(i,j) \ne (i',j')$} \,.
    \end{displaymath}
    Therefore
    \begin{multline*}
        |\hat\chi(A) - \hat\chi(A')|^2 = \textsum{i,j}{} (\chi(a_i) - \chi(a'_j))^2 |\pi_i \circ
        \pi'_j|^2 \\
        \le \Lip(\chi)^2 \textsum{i,j}{} (a_i - a'_j)^2 |\pi_i \circ \pi'_j|^2 =
        \Lip(\chi)^2 |A - A'|^2 \,. \qedhere
    \end{multline*}
\end{proof}

\begin{lemma}
    \label{lem:trace}
    Suppose $g \in \textstyle\bigodot^2(\R^{\adim},\R)$, $g(v,v) \ge 0$ for $v \in \R^{\adim}$, $Q
    \in \SymEnd{\R^{\adim}}$, and $\|Q\| \le 1$. Then
    \begin{displaymath}
        \textsum{j=1}{\adim} g(Q(\varepsilon_j),Q(\varepsilon_j)) \le \textsum{j=1}{\adim}
        g(\varepsilon_j,\varepsilon_j) \,.
    \end{displaymath}
\end{lemma}

\begin{proof}
    Both sides are traces of symmetric bilinear forms with respect to~$\bullet$; hence they do not
    depend on the orthonormal basis~\cite[1.7.10]{Federer1969}. Choose an~orthonormal basis $u_1,
    \ldots, u_{\adim}$ with $Q(u_j) = q_j u_j$~\cite[1.7.3]{Federer1969}. Then $|q_j| \le 1$ and
    $\textsum{j}{} g(Q(u_j),Q(u_j)) = \textsum{j}{} q_j^2 g(u_j,u_j) \le \textsum{j}{}
    g(u_j,u_j)$.
\end{proof}

\begin{lemma}
    \label{lem:hull}
    If $Z$ is a~finite dimensional normed space and $A \subseteq Z$ is compact, then $\conv A$ is
    compact.
\end{lemma}

\begin{proof}
    Let $d = \dim Z$. Suppose $z = \textsum{i=0}{m} c_i a_i$ with $a_i \in A$, $c_i > 0$,
    $\textsum{i}{} c_i = 1$, and $m > d$. The vectors $a_1 - a_0, \ldots, a_m - a_0$ are linearly
    dependent; hence there are $b_0, \ldots, b_m \in \R$, not all zero, with $\textsum{i}{} b_i
    a_i = 0$ and $\textsum{i}{} b_i = 0$. Some $b_i$ is positive. With $\tau = \min \{ c_i / b_i :
    b_i > 0 \}$ we have $z = \textsum{i}{} (c_i - \tau b_i) a_i$, where $c_i - \tau b_i \ge 0$,
    $\textsum{i}{} (c_i - \tau b_i) = 1$, and at least one coefficient vanishes. By induction on
    $m$, $\conv A$ is the image of the compact set $A^{d+1} \times (\R^{d+1} \cap \{ c : c_i \ge
    0,\ \textsum{i}{} c_i = 1 \})$ under the continuous map $[(a,c) \mapsto \textsum{i}{} c_i
    a_i]$.
\end{proof}

\begin{lemma}
    \label{lem:barycentre}
    Suppose $Z$ is a~finite dimensional normed space, $K \subseteq Z$ is compact and convex, $\nu$
    is a~measure over a~set~$\Omega$ with $0 < \nu(\Omega) < \infty$, and $f : \Omega \to K$ is
    $\nu$~measurable. Then $\nu(\Omega)^{-1} \textint{}{} f \ud \nu \in K$.
\end{lemma}

\begin{proof}
    The integral exists by~\cite[2.4.12]{Federer1969}, because $f$ is bounded. The membership is
    part of Jensen's inequality~\cite[2.4.19]{Federer1969}, applied to the closed convex set~$K$
    and the zero function on~$K$.
\end{proof}

\section{Heat flow of finite measures with compact support}
\label{sec:heat}

\begin{miniremark}
    \label{mr:sec_heat}
    This section collects the properties of the heat kernel and of heat extensions of finite
    measures with compact support used in sections~\ref{sec:estimates} and~\ref{sec:entropy}. It
    introduces the classes of functions of polynomial growth and of Gaussian decay, which justify
    all integrations by parts on~$\R^{\adim}$ and all differentiations under the integral sign.
    It ends with the log-semiconvexity of heat extensions~\ref{lem:semiconvex}
    (\cite[(3.16)]{DePhilippisGennaioliPigatiRindler2026}) and the integration by parts formula
    \ref{lem:ibp} (\cite[(3.17)]{DePhilippisGennaioliPigatiRindler2026}).
\end{miniremark}

\begin{definition}
    \label{def:heat}
    For $0 < s < \infty$ and $z \in \R^{\adim}$ we set $p_s(z) = (4 \pi s)^{-\adim/2} \exp(-|z|^2
    / (4s))$. Suppose $\nu$ is a~Borel measure over~$\R^{\adim}$, $0 < r < \infty$, $\nu(\R^{\adim})
    < \infty$, $\nu(\R^{\adim} \without \cball{0}{r}) = 0$, $Z$ is a~finite dimensional normed
    space, and $f : \R^{\adim} \to Z$ is bounded and $\nu$~measurable. The \emph{heat extension}
    of $(\nu,f)$ is the map $u : (0,\infty) \times \R^{\adim} \to Z$ given by
    \begin{displaymath}
        u(s,x) = \textint{}{} p_s(x-y) f(y) \ud \nu y
        \quad \text{for $0 < s < \infty$ and $x \in \R^{\adim}$} \,,
    \end{displaymath}
    where the integral is that of~\cite[2.4.12]{Federer1969}.
\end{definition}

\begin{lemma}
    \label{lem:kernel}
    For $0 < s < \infty$, $z, v, w \in \R^{\adim}$, and $j \in \nat$ the following six statements
    hold.
    \begin{enumerate}
        \item \label{i:k:mass} $\textint{}{} p_s \ud \LM^{\adim} = 1$.
        \item \label{i:k:moment} $\textint{}{} |z|^2 p_s(z) \ud \LM^{\adim} z = 2 \adim s$.
        \item \label{i:k:heat} $\partial_s p_s(z) = \Delta p_s(z)$.
        \item \label{i:k:bound} $\|\uD^j p_s(z)\| \le c(\adim,j) \, s^{-j/2} (1 + |z| s^{-1/2})^j \,
        p_s(z)$, where $0 < c(\adim,j) < \infty$ depends only on $\adim$ and~$j$.
        \item \label{i:k:formulas} $\uD p_s(z)(v) = - (2s)^{-1} (z \bullet v) \, p_s(z)$ and
        $\uD^2 p_s(z)(v,w) = \bigl( (4s^2)^{-1} (z \bullet v)(z \bullet w) - (2s)^{-1} v \bullet w
        \bigr) p_s(z)$.
        \item \label{i:k:max} $p_s(z) \le (4 \pi s)^{-\adim/2}$.
    \end{enumerate}
\end{lemma}

\begin{proof}
    Statements~\ref{i:k:formulas} and~\ref{i:k:max} follow by direct computation,
    and~\ref{i:k:heat} follows from~\ref{i:k:formulas}, since both sides equal $((4s^2)^{-1}|z|^2
    - \adim (2s)^{-1}) p_s(z)$. By Fubini's theorem~\cite[2.6.2]{Federer1969}, \ref{i:k:mass}
    and~\ref{i:k:moment} reduce to the one-dimensional identities $\textint{}{} \exp(-\tau^2/(4s))
    \ud \LM^1 \tau = (4 \pi s)^{1/2}$ and $\textint{}{} \tau^2 \exp(-\tau^2/(4s)) \ud \LM^1 \tau =
    2s (4\pi s)^{1/2}$. For~\ref{i:k:bound}, note $p_s(z) = s^{-\adim/2} p_1(s^{-1/2} z)$. By
    induction on~$j$, $\uD^j p_1(w) = P_j(w) p_1(w)$ for a~polynomial function $P_j : \R^{\adim}
    \to \textstyle\bigodot^j(\R^{\adim},\R)$ of degree at most~$j$; indeed $P_0 = 1$ and
    \begin{displaymath}
        P_{j+1}(w)(v_0, \ldots, v_j) = \uD P_j(w)(v_0)(v_1, \ldots, v_j) - 2^{-1} (w \bullet v_0)
        P_j(w)(v_1, \ldots, v_j) \,.
    \end{displaymath}
    Hence $\|P_j(w)\| \le c(\adim,j)(1 + |w|)^j$ and $\uD^j p_s(z) =
    s^{-\adim/2} s^{-j/2} \uD^j p_1(s^{-1/2} z)$, which gives~\ref{i:k:bound}.
\end{proof}

\begin{lemma}
    \label{lem:heat_extension}
    Suppose $\nu$, $r$, $Z$, $f$ and~$u$ are as in~\ref{def:heat}, and $c(\adim,j)$ is as
    in~\ref{lem:kernel}\ref{i:k:bound}. Then the following four statements hold.
    \begin{enumerate}
        \item \label{i:he:smooth} The map~$u$ is of class~$\infty$, $\partial_s u = \Delta u$, and
        $\uD^j u_s(x) = \textint{}{} \uD^j p_s(x-y) f(y) \ud \nu y$ for $0 < s < \infty$, $x \in
        \R^{\adim}$, and $j \in \nat$.
        \item \label{i:he:bound} For $0 < s < \infty$, $x \in \R^{\adim}$, and $j \in \nat$,
        \begin{displaymath}
            \|\uD^j u_s(x)\| \le c(\adim,j) \, s^{-j/2} \bigl( 1 + (|x| + r) s^{-1/2} \bigr)^j
            \textint{}{} p_s(x-y) |f(y)| \ud \nu y \,.
        \end{displaymath}
        \item \label{i:he:gauss} For $0 < s < \infty$ and $x \in \R^{\adim}$,
        \begin{displaymath}
            \textint{}{} p_s(x-y) |f(y)| \ud \nu y \le (4 \pi s)^{-\adim/2} \exp \bigl(
            - \max \{ |x| - r, 0 \}^2 / (4s) \bigr) \textint{}{} |f| \ud \nu \,;
        \end{displaymath}
        if moreover $Z = \R$ and $f \ge 0$, then $u(s,x) \ge (4 \pi s)^{-\adim/2} \exp(-(|x| +
        r)^2/(4s)) \textint{}{} f \ud \nu$.
        \item \label{i:he:mass} $\textint{}{} u_s \ud \LM^{\adim} = \textint{}{} f \ud \nu$ for $0
        < s < \infty$.
    \end{enumerate}
\end{lemma}

\begin{proof}
    Suppose $0 < s_0 < s_1 < \infty$. By~\ref{lem:kernel}\ref{i:k:bound}\ref{i:k:heat}, every
    derivative of $[(s,z) \mapsto p_s(z)]$ is bounded on $[s_0,s_1] \times \R^{\adim}$. Hence
    differentiation under the integral sign, justified by the mean value theorem and Lebesgue's
    bounded convergence theorem~\cite[2.4.9]{Federer1969}, gives~\ref{i:he:smooth}, where $\partial_s
    u = \Delta u$ follows from~\ref{lem:kernel}\ref{i:k:heat}. Since $|x - y| \le |x| + r$ for
    $y \in \cball{0}{r}$, \ref{i:he:bound} follows from~\ref{i:he:smooth}
    and~\ref{lem:kernel}\ref{i:k:bound}. Since $\max \{ |x| - r, 0 \} \le |x - y| \le |x| + r$ for $y
    \in \cball{0}{r}$, we obtain~\ref{i:he:gauss}. Finally, \ref{i:he:mass} follows from Fubini's
    theorem~\cite[2.6.2]{Federer1969} and~\ref{lem:kernel}\ref{i:k:mass}.
\end{proof}

\begin{definition}
    \label{def:classes}
    Suppose $Z$ is a~finite dimensional normed space. A~map $f : \R^{\adim} \to Z$ of
    class~$\infty$ has \emph{polynomial growth} if for each $j \in \nat$ there is $0 < c <
    \infty$ with $\|\uD^j f(x)\| \le c (1 + |x|)^c$ for $x \in \R^{\adim}$; it has
    \emph{Gaussian decay} if for each $j \in \nat$ there is $0 < c < \infty$ with $\|\uD^j f(x)\|
    \le c \exp(-|x|^2/c)$ for $x \in \R^{\adim}$.

    Suppose $0 < T < \infty$. A~map $f : \Omega_T \to Z$ of class~$\infty$ has \emph{polynomial
    growth before~$T$} if for each $j \in \nat$ and each compact interval $J \subseteq \R \cap \{ t
    : t < T \}$ there is $0 < c < \infty$ with $\|\uD^j f(t,x)\| \le c (1 + |x|)^c$ for $t \in
    J$ and $x \in \R^{\adim}$; \emph{Gaussian decay before~$T$} is defined likewise with $c
    \exp(-|x|^2/c)$ in place of $c (1 + |x|)^c$.

    A~locally Lipschitzian map $e : \R^{\adim} \to \R^{\adim}$ has \emph{polynomial growth} if
    there is $0 < c < \infty$ with $|e(x)| + \Lip(e|\cball{x}{1}) \le c (1 + |x|)^c$ for $x \in
    \R^{\adim}$.
\end{definition}

\begin{lemma}
    \label{lem:classes}
    Suppose $0 < T < \infty$, $Z_1$, $Z_2$, $Z_3$ are finite dimensional normed spaces, and $\beta :
    Z_1 \times Z_2 \to Z_3$ is bilinear. Then the following three statements hold.
    \begin{enumerate}
        \item \label{i:cl:product} If $f_1 : \Omega_T \to Z_1$ and $f_2 : \Omega_T \to Z_2$ have
        polynomial growth before~$T$, then so does $[(t,x) \mapsto \beta(f_1(t,x),f_2(t,x))]$;
        if moreover $f_2$ has Gaussian decay before~$T$, then so does this map. The same holds
        for maps on~$\R^{\adim}$ and the classes without time.
        \item \label{i:cl:dt} If $f : \Omega_T \to Z_1$ has Gaussian decay before~$T$, then the
        map $[t \mapsto \textint{}{} f_t \ud \LM^{\adim}]$ is of class~$1$ on $\R \cap \{ t : t < T
        \}$, and its derivative at~$t$ is $\textint{}{} \partial_t f(t,x) \ud \LM^{\adim} x$.
        \item \label{i:cl:div} If $g : \R^{\adim} \to \R^{\adim}$ is of class~$1$ and there is $0 <
        c < \infty$ with $|g(x)| + \|\uD g(x)\| \le c \exp(-|x|^2/c)$ for $x \in \R^{\adim}$,
        then $\textint{}{} \mathrm{div}\, g \ud \LM^{\adim} = 0$. If $f : \R^{\adim} \to \R$ has
        Gaussian decay, then $\textint{}{} \Delta f \ud \LM^{\adim} = 0$.
    \end{enumerate}
\end{lemma}

\begin{proof}
    Statement~\ref{i:cl:product} follows from the Leibniz formula~\cite[3.1.11]{Federer1969} for
    $\uD^j \beta(f_1,f_2)$ and from $(1 + |x|)^c \exp(-|x|^2/c') \le c'' \exp(-|x|^2/(2c'))$,
    where $c''$ depends on $c$ and~$c'$. For~\ref{i:cl:dt}, the difference quotients in~$t$ of
    $f(\cdot,x)$ are bounded on each compact interval by $c \exp(-|x|^2/c)$, by the mean value
    theorem; Lebesgue's bounded convergence theorem~\cite[2.4.9]{Federer1969} gives the derivative
    and its continuity. For~\ref{i:cl:div}, $\textint{}{} \mathrm{div}\, g \ud \LM^{\adim} =
    \textsum{j=1}{\adim} \textint{}{} \uD g(x)(\varepsilon_j) \bullet \varepsilon_j \ud
    \LM^{\adim} x$. By Fubini's theorem~\cite[2.6.2]{Federer1969}, the $j$-th term is an~iterated
    integral whose inner integral is the integral over a~line parallel to~$\varepsilon_j$ of the
    derivative of $[x \mapsto g(x) \bullet \varepsilon_j]$ along that line; it vanishes, since the
    function tends to zero at infinity. Finally, $\Delta f = \mathrm{div} \grad f$.
\end{proof}

\begin{lemma}
    \label{lem:semiconvex}
    Suppose $\nu$ and~$r$ are as in~\ref{def:heat}, $\nu(\R^{\adim}) > 0$, and $u$ is the heat
    extension of $(\nu,1)$. Then, for $0 < s < \infty$ and $x, v \in \R^{\adim}$,
    \begin{displaymath}
        |\grad \log u_s(x)| \le (2s)^{-1} (|x| + r) \,, \qquad
        \uD^2 \log u_s(x)(v,v) \ge - (2s)^{-1} |v|^2 \,.
    \end{displaymath}
\end{lemma}

\begin{proof}
    By~\ref{lem:heat_extension}\ref{i:he:gauss}, $u_s(x) > 0$. Let $\pi = u_s(x)^{-1} \nu
    \restrict [y \mapsto p_s(x-y)]$; then $\pi(\R^{\adim}) = 1$ and $\pi(\R^{\adim} \without
    \cball{0}{r}) = 0$. Let $w(y) = (x - y) \bullet v$ for $y \in \R^{\adim}$.
    By~\ref{lem:heat_extension}\ref{i:he:smooth} and~\ref{lem:kernel}\ref{i:k:formulas},
    \begin{displaymath}
        \frac{\uD u_s(x)(v)}{u_s(x)} = - \frac{1}{2s} \textint{}{} w \ud \pi \,, \qquad
        \frac{\uD^2 u_s(x)(v,v)}{u_s(x)} = \frac{1}{4s^2} \textint{}{} w^2 \ud \pi -
        \frac{|v|^2}{2s} \,.
    \end{displaymath}
    Since $|w(y)| \le (|x| + r) |v|$ for $y \in \cball{0}{r}$, the first inequality follows. Since
    $\uD^2 \log u_s(x)(v,v) = u_s(x)^{-1} \uD^2 u_s(x)(v,v) - (u_s(x)^{-1} \uD u_s(x)(v))^2$ and
    $(\textint{}{} w \ud \pi)^2 \le \textint{}{} w^2 \ud \pi$ by H\"older's
    inequality~\cite[2.4.14]{Federer1969}, the second inequality follows.
\end{proof}

\begin{lemma}
    \label{lem:ibp}
    Suppose $\rho : \R^{\adim} \to \R$ is of class~$\infty$, $\rho(x) > 0$ for $x \in
    \R^{\adim}$, $\rho$ has Gaussian decay, $\uD \log \rho$ has polynomial growth, and $e :
    \R^{\adim} \to \R^{\adim}$ is a~locally Lipschitzian map of polynomial growth. Then $e$ is
    differentiable at $\LM^{\adim}$~almost all points and
    \begin{displaymath}
        - \textint{}{} \rho \, \uD^2 \log \rho(e,e) \ud \LM^{\adim}
        = \textint{}{} \rho \bigl[ (\mathrm{div}\, e + \uD \log \rho(e))^2 - \trace(\uD e \circ
        \uD e) \bigr] \ud \LM^{\adim} \,.
    \end{displaymath}
\end{lemma}

\begin{proof}
    Differentiability almost everywhere holds by Rademacher's theorem~\cite[3.1.6]{Federer1969}.
    Let $h = \log \rho$. Suppose first that $e$ is of class~$\infty$ and $e$, $\uD e$, $\uD^2 e$
    have polynomial growth in the sense that their norms are bounded by $c(1 + |x|)^c$. Set $b =
    \mathrm{div}\, e + \uD h(e)$. Since $\uD \rho = \rho \uD h$, we have $\mathrm{div}(\rho e) =
    \rho b$. Since $\uD^2 e(x)$ is symmetric,
    \begin{displaymath}
        \mathrm{div} \bigl[ x \mapsto \uD e(x)(e(x)) \bigr] = \uD(\mathrm{div}\, e)(e) + \trace(\uD
        e \circ \uD e) \,.
    \end{displaymath}
    The maps $b \rho e$ and $[x \mapsto \rho(x) \uD e(x)(e(x))]$ satisfy the hypothesis
    of~\ref{lem:classes}\ref{i:cl:div} by the Leibniz formula~\cite[3.1.11]{Federer1969}, because
    $\rho$ has Gaussian decay and $e$, $\uD e$, $\uD^2 e$, $\uD h$, $\uD^2 h$ have polynomial
    growth. Since
    $\mathrm{div}(b \rho e) = \rho b^2 + \rho \uD b(e)$ and $\uD b(e) = \uD(\mathrm{div}\, e)(e) +
    \uD^2 h(e,e) + \uD h(\uD e(e))$, we obtain
    \begin{align}
        0 & = \textint{}{} \rho b^2 \ud \LM^{\adim} + \textint{}{} \rho \, \uD(\mathrm{div}\, e)(e)
        \ud \LM^{\adim} + \textint{}{} \rho \, \uD^2 h(e,e) \ud \LM^{\adim} + \textint{}{} \rho
        \, \uD h(\uD e(e)) \ud \LM^{\adim} \,, \label{eq:ibp:1} \\
        0 & = \textint{}{} \rho \, \uD h(\uD e(e)) \ud \LM^{\adim} + \textint{}{} \rho \,
        \uD(\mathrm{div}\, e)(e) \ud \LM^{\adim} + \textint{}{} \rho \trace(\uD e \circ \uD e) \ud
        \LM^{\adim} \,, \label{eq:ibp:2}
    \end{align}
    where~\eqref{eq:ibp:2} expresses $\textint{}{} \mathrm{div}[x \mapsto \rho(x) \uD
    e(x)(e(x))] \ud \LM^{\adim} = 0$. Subtracting~\eqref{eq:ibp:2} from~\eqref{eq:ibp:1} gives the
    conclusion.

    In the general case let $\Phi \in \dspace{\R^{\adim}}{\R}$ satisfy $\Phi \ge 0$, $\spt \Phi
    \subseteq \cball{0}{1}$ and $\textint{}{} \Phi \ud \LM^{\adim} = 1$, let $\Phi_\epsilon(x) =
    \epsilon^{-\adim} \Phi(\epsilon^{-1} x)$ as in~\cite[4.1.2]{Federer1969}, and let
    $e_\epsilon(x) = \textint{}{} \Phi_\epsilon(x-y) e(y) \ud \LM^{\adim} y$ for $0 < \epsilon \le
    1$. Then $e_\epsilon$ is of class~$\infty$. For $x, v \in \R^{\adim}$, the difference
    quotients $\tau^{-1}(e(y + \tau v) - e(y))$ are bounded for $y \in \cball{x}{1}$ and $0 <
    |\tau| \le 1$, and converge to $\uD e(y)(v)$ for $\LM^{\adim}$~almost all~$y$; hence, by
    Lebesgue's bounded convergence theorem~\cite[2.4.9]{Federer1969},
    \begin{displaymath}
        \uD e_\epsilon(x) = \textint{}{} \Phi_\epsilon(x-y) \uD e(y) \ud \LM^{\adim} y \,.
    \end{displaymath}
    With $c$ as in~\ref{def:classes} for~$e$, we have $\|\uD e(y)\| \le c(1 + |y|)^c$ at points of
    differentiability, and therefore
    \begin{displaymath}
        |e_\epsilon(x)| + |\uD e_\epsilon(x)| \le (1 + \adim^{1/2}) \, c \, (2 + |x|)^c
        \quad \text{for $x \in \R^{\adim}$ and $0 < \epsilon \le 1$} \,.
    \end{displaymath}
    Moreover, $\uD^2 e_\epsilon(x)(v,w) = \textint{}{} \uD \Phi_\epsilon(x-y)(w) \, \uD e(y)(v) \ud
    \LM^{\adim} y$, so the first case applies to~$e_\epsilon$. As $\epsilon \to 0+$,
    $e_\epsilon(x) \to e(x)$ for every~$x$, and $\uD e_\epsilon(x) \to \uD e(x)$ at every Lebesgue
    point~$x$ of~$\uD e$, hence for $\LM^{\adim}$~almost all~$x$ by~\cite[2.9.9]{Federer1969}. The
    integrands of the conclusion for~$e_\epsilon$ are bounded by $\rho(x) \, c' (1 + |x|)^{c'}$,
    where $c'$ does not depend on~$\epsilon$, and this function is $\LM^{\adim}$~summable.
    Lebesgue's bounded convergence theorem~\cite[2.4.9]{Federer1969} yields the conclusion
    for~$e$.
\end{proof}

\section{Heat estimates}
\label{sec:estimates}

\begin{miniremark}
    \label{mr:standing}
    This section reproduces~\cite[section~3]{DePhilippisGennaioliPigatiRindler2026}: the heated
    objects~\ref{def:heat_objects}, their identities~\ref{lem:identities}, their
    regularity~\ref{lem:regularity}, the calculus of the heat operators~\ref{lem:heat_calculus},
    and the two integral estimates~\ref{lem:DGPR31} and~\ref{lem:DGPR32}.

    In this section and in section~\ref{sec:entropy} we assume the following hypotheses, called
    \emph{(S)}. The space $Y$ is a~finite dimensional inner product space, $K \subseteq
    \Hom(\R^{\adim},Y)$ is nonempty, compact and convex, $0 < \lambda < \infty$, $0 < \Lambda <
    \infty$, and
    \begin{displaymath}
        \sv{\vdim}(X) \ge \lambda \quad \text{and} \quad |X| \le \Lambda
        \qquad \text{for $X \in K$} \,.
    \end{displaymath}
    The measure $\mu$ is a~Borel measure over~$\R^{\adim}$ with $\mu(\R^{\adim}) = 1$ and
    $\mu(\R^{\adim} \without \cball{0}{r_0}) = 0$ for some $0 < r_0 < \infty$. The map $X :
    \R^{\adim} \to K$ is $\mu$~measurable, and $\sigma : \ccspace{\R^{\adim},\Hom(\R^{\adim},Y)}
    \to \R$ is given by
    \begin{displaymath}
        \sigma(\Phi) = \textint{}{} X(x) \bullet \Phi(x) \ud \mu x
        \quad \text{for $\Phi \in \ccspace{\R^{\adim},\Hom(\R^{\adim},Y)}$} \,.
    \end{displaymath}
    The distribution $\Upsilon \in \distr{\R^{\adim}}{Y}$ is representable by integration,
    $\delta = \|\Upsilon\|(\R^{\adim}) < \infty$, and
    \begin{displaymath}
        \Upsilon(g) = - \sigma(\uD g) \quad \text{for $g \in \dspace{\R^{\adim}}{Y}$} \,.
    \end{displaymath}
    Finally, $0 < T < \infty$, $\mu_0$ is a~Borel measure over~$\R^{\adim}$ with $\mu_0(A) \le
    \mu(A)$ for $A \subseteq \R^{\adim}$, and $\gamma = \mu_0(\R^{\adim}) > 0$. Thus $0 < \gamma
    \le 1$. The map $\dUps : \R^{\adim} \to \Hom(Y,\R)$ is that of~\ref{mr:functionals},
    modified on a~$\|\Upsilon\|$ null set so that $\|\dUps(z)\| \le 1$ for $z \in \R^{\adim}$.

    In the language of~\cite{DePhilippisGennaioliPigatiRindler2026}, $\sigma$ is the measure $X
    \mu$, $\Upsilon$ is its row-wise divergence, $\|\Upsilon\|$ is $|\mathrm{div}\, \sigma|$, and
    the hypotheses $\mu(\R^{\adim}) = 1$ and $\mu(\R^{\adim} \without \cball{0}{r_0}) = 0$ are
    those at the beginning of~\cite[section~3]{DePhilippisGennaioliPigatiRindler2026}.
\end{miniremark}

\begin{definition}
    \label{def:heat_objects}
    Assume (S). For $(t,x) \in \Omega_T$ we define
    \begin{gather}
        U(t,x) = \textint{}{} p_{T-t}(x-y) \ud \mu y \,, \qquad
        W(t,x) = \textint{}{} p_{T-t}(x-y) X(y) \ud \mu y \,, \\
        R(t,x) = \textint{}{} p_{T-t}(x-y) \ud \|\Upsilon\| y \,, \qquad
        \Psi(t,x) = \textint{}{} p_{T-t}(x-y) \dUps(y) \ud \|\Upsilon\| y \,, \\
        S(t,x) = U(t,x) + (T-t)^{1/2} R(t,x) \,, \qquad
        \rho(t,x) = \textint{}{} p_{T-t}(x-y) \ud \mu_0 y \,.
    \end{gather}
    Thus $W(t,x) \in \Hom(\R^{\adim},Y)$ and $\Psi(t,x) \in \Hom(Y,\R)$. Whenever $S(t,x) > 0$ and
    $\rho(t,x) > 0$ we further define
    \begin{gather}
        B(t,x) = S(t,x)^{-1} W(t,x) \,, \qquad
        \zeta(t,x) = \grad \log S_t(x) \,, \qquad
        \eta(t,x) = S(t,x)^{-1} \Psi(t,x) \,, \\
        \kappa(t,x) = \frac{R(t,x)}{2 (T-t)^{1/2} S(t,x)} \,, \qquad
        q(t,x) = \frac{\rho(t,x)}{S(t,x)} \,, \qquad
        I(\rho_t) = \textint{}{} \rho_t |\grad \log \rho_t|^2 \ud \LM^{\adim} \,.
    \end{gather}
    For a~map $f : \Omega_T \to Z$ of class~$2$ into a~finite dimensional normed space~$Z$ we
    define $\heatop f, \heatop_\zeta f : \Omega_T \to Z$ by
    \begin{displaymath}
        \heatop f(t,x) = \partial_t f(t,x) + \Delta f_t(x) \,, \qquad
        \heatop_\zeta f(t,x) = \heatop f(t,x) + 2 \uD f_t(x)(\zeta(t,x)) \,.
    \end{displaymath}
    These are $P_{T-t}\mu$, $P_{T-t}\sigma$, $P_{T-t}|\mathrm{div}\,\sigma|$,
    $P_{T-t}\mathrm{div}\,\sigma$, $\mathcal{H}$ and $\mathcal{H}_\zeta$
    of~\cite[(3.1), (3.2), (3.4), (3.6)]{DePhilippisGennaioliPigatiRindler2026}, with $W$ in place
    of~$V$; positivity of $S$ and~$\rho$ is proved in~\ref{lem:identities}.
\end{definition}

\begin{lemma}
    \label{lem:identities}
    Assume (S). Then the following seven statements hold.
    \begin{enumerate}
        \item \label{i:id:spt} $\|\Upsilon\|(\R^{\adim} \without \cball{0}{r_0}) = 0$.
        \item \label{i:id:div} The maps $U$, $W$, $R$, $\Psi$, $S$, $\rho$ are of class~$\infty$,
        and $\mathrm{div}\, W_t(x) = \Psi(t,x)$ and $\|\Psi(t,x)\| \le R(t,x)$ for $(t,x) \in
        \Omega_T$.
        \item \label{i:id:heat} $\heatop U = 0$, $\heatop W = 0$, $\heatop R = 0$, $\heatop \Psi =
        0$, $\heatop \rho = 0$, and $\heatop S = - \kappa S$ on~$\Omega_T$.
        \item \label{i:id:mass} $\textint{}{} U_t \ud \LM^{\adim} = 1$, $\textint{}{} R_t \ud
        \LM^{\adim} = \delta$, $\textint{}{} S_t \ud \LM^{\adim} = 1 + (T-t)^{1/2} \delta$, and
        $\textint{}{} \rho_t \ud \LM^{\adim} = \gamma$ for $t < T$.
        \item \label{i:id:bounds} On~$\Omega_T$ we have $0 < \rho \le U \le S$, $\kappa \ge 0$,
        $U^{-1} W \in K$, $|B| \le \Lambda$, and $\|\eta\| \le (T-t)^{-1/2}$.
        \item \label{i:id:divB} For $(t,x) \in \Omega_T$ and $y \in Y$,
        \begin{displaymath}
            \langle y, \mathrm{div}\, B_t(x) \rangle + B(t,x)(\zeta(t,x)) \bullet y = \langle y,
            \eta(t,x) \rangle \,.
        \end{displaymath}
        \item \label{i:id:sv} If $(t,x) \in \Omega_T$ and $U(t,x) \ge (T-t)^{1/2} R(t,x)$, then
        $\sv{\vdim}(B(t,x)) \ge \lambda/2$.
    \end{enumerate}
\end{lemma}

\begin{proof}
    If $g \in \dspace{\R^{\adim}}{Y}$ and $\spt g \cap \cball{0}{r_0} = \varnothing$, then $\uD
    g(x) = 0$ for $\mu$~almost all~$x$, hence $\Upsilon(g) = -\sigma(\uD g) = 0$;
    thus~\ref{i:id:spt} follows from~\ref{mr:functionals}. By~\ref{i:id:spt}, $U$, $W$, $R$,
    $\Psi$ and $\rho$ are of the form $[(t,x) \mapsto u(T-t,x)]$, where $u$ is the heat extension
    of $(\mu,1)$, $(\mu,X)$, $(\|\Upsilon\|,1)$, $(\|\Upsilon\|,\dUps)$, $(\mu_0,1)$, all with
    $r = r_0$. Hence they and~$S$ are of class~$\infty$ by~\ref{lem:heat_extension}\ref{i:he:smooth}.

    To prove $\mathrm{div}\, W_t(x) = \Psi(t,x)$, fix $(t,x) \in \Omega_T$ and $y \in Y$, write $s
    = T - t$, and let $g(z) = p_s(x-z) y$ for $z \in \R^{\adim}$. Choose $\varphi \in
    \dspace{\R^{\adim}}{\R}$ with $\varphi(z) = 1$ for $z \in \oball{0}{r_0 + 1}$. Then $\varphi g
    \in \dspace{\R^{\adim}}{Y}$. By~\ref{i:id:spt},
    \begin{displaymath}
        \Upsilon(\varphi g) = \textint{}{} p_s(x-z) \langle y, \dUps(z) \rangle \ud \|\Upsilon\| z
        = \langle y, \Psi(t,x) \rangle \,.
    \end{displaymath}
    Since $\uD(\varphi g)(z) = \uD g(z) = - \uD p_s(x-z) \cdot y$ for $z \in \oball{0}{r_0 + 1}$,
    \begin{displaymath}
        \Upsilon(\varphi g) = - \sigma(\uD(\varphi g)) = \textint{}{} X(z) \bullet (\uD p_s(x-z)
        \cdot y) \ud \mu z = \textint{}{} \uD p_s(x-z)(\transpose{X(z)}(y)) \ud \mu z \,.
    \end{displaymath}
    On the other hand, $\uD W_t(x)(v) = \textint{}{} \uD p_s(x-z)(v) X(z) \ud \mu z$
    by~\ref{lem:heat_extension}\ref{i:he:smooth}, hence
    \begin{displaymath}
        \langle y, \mathrm{div}\, W_t(x) \rangle = \textsum{j=1}{\adim} \textint{}{} \uD
        p_s(x-z)(\varepsilon_j) \, (X(z)(\varepsilon_j) \bullet y) \ud \mu z = \textint{}{} \uD
        p_s(x-z)(\transpose{X(z)}(y)) \ud \mu z \,.
    \end{displaymath}
    Combining the three displays gives $\mathrm{div}\, W_t(x) = \Psi(t,x)$. Since $\|\dUps(z)\| = 1$
    for $\|\Upsilon\|$~almost all~$z$, we have $\|\Psi(t,x)\| \le R(t,x)$. This
    proves~\ref{i:id:div}.

    Since $\partial_t [u(T-t,x)] = - \partial_s u(T-t,x) = - \Delta u_{T-t}(x)$
    by~\ref{lem:heat_extension}\ref{i:he:smooth}, the first five equations of~\ref{i:id:heat} hold.
    Moreover, $\heatop((T-t)^{1/2} R) = -2^{-1} (T-t)^{-1/2} R + (T-t)^{1/2} \heatop R = - \kappa
    S$, whence $\heatop S = -\kappa S$. Statement~\ref{i:id:mass} follows
    from~\ref{lem:heat_extension}\ref{i:he:mass}.

    By~\ref{lem:heat_extension}\ref{i:he:gauss}, $U > 0$ and $\rho > 0$; since $\mu_0 \le \mu$, we
    have $\rho \le U$, and $U \le S$ by definition. Hence $B$, $\zeta$, $\eta$, $\kappa$, $q$ are
    defined on~$\Omega_T$, and $\kappa \ge 0$. By~\ref{lem:barycentre} applied to the measure $\mu
    \restrict [y \mapsto p_{T-t}(x-y)]$ and the map~$X$, we have $U(t,x)^{-1} W(t,x) \in K$; hence
    $|B| = (U/S) \, |U^{-1} W| \le \Lambda$. Since $(T-t)^{1/2} R \le S$ and $\|\Psi\| \le R$, we
    obtain $\|\eta\| \le S^{-1} R \le (T-t)^{-1/2}$. This proves~\ref{i:id:bounds}.

    Since $\uD B_t(x)(v) = S^{-1} \uD W_t(x)(v) - S^{-2} \uD S_t(x)(v) \, W(t,x)$ and
    $\textsum{j}{} \uD S_t(x)(\varepsilon_j) W(t,x)(\varepsilon_j) = W(t,x)(\grad S_t(x))$,
    \begin{displaymath}
        \langle y, \mathrm{div}\, B_t(x) \rangle = S^{-1} \langle y, \mathrm{div}\, W_t(x) \rangle
        - S^{-1} B(t,x)(\grad S_t(x)) \bullet y \,,
    \end{displaymath}
    with $S = S(t,x)$; together with~\ref{i:id:div} this gives~\ref{i:id:divB}.

    Suppose $U(t,x) \ge (T-t)^{1/2} R(t,x)$. Then $S(t,x) \le 2 U(t,x)$. Moreover, $B(t,x) = c
    \, U(t,x)^{-1} W(t,x)$, where $c = U(t,x) / S(t,x) \ge 1/2$ and $U(t,x)^{-1} W(t,x) \in K$.
    Hence $\sv{\vdim}(B(t,x)) \ge \lambda/2$ by~\ref{lem:sv}\ref{i:sv:hom}. This proves~\ref{i:id:sv}.
\end{proof}

\begin{lemma}
    \label{lem:regularity}
    Assume (S). Then $U$, $W$, $R$, $\Psi$, $S$ and $\rho$ have Gaussian decay before~$T$, and
    $\log S$, $\log \rho$, $q$, $\log q$, $B$, $\zeta$, $\eta$ and $\kappa$ are of class~$\infty$
    and have polynomial growth before~$T$. In particular, for $t < T$, the function $\rho_t$
    satisfies the hypotheses on~$\rho$ in~\ref{lem:ibp}.
\end{lemma}

\begin{proof}
    Let $J \subseteq \R \cap \{ t : t < T \}$ be a~compact interval; then $T - t$ ranges over
    a~compact subset of $(0,\infty)$ for $t \in J$. By~\ref{lem:heat_extension}\ref{i:he:smooth},
    each derivative of order~$j$ with respect to $(t,x)$ of $U$, $W$, $R$, $\Psi$, $\rho$ is
    a~linear combination of derivatives with respect to~$x$ of order at most~$2j$; the same holds
    for~$S$, since $[t \mapsto (T-t)^{1/2}]$ is smooth on~$J$. By~\ref{lem:heat_extension}
    \ref{i:he:bound}\ref{i:he:gauss} and $\max \{ |x| - r_0, 0 \}^2 \ge 2^{-1} |x|^2 - r_0^2$,
    these maps have Gaussian decay before~$T$. Moreover, \ref{lem:heat_extension}\ref{i:he:bound}
    together with $\textint{}{} p_{T-t}(x-y) |X(y)| \ud \mu y \le \Lambda U(t,x)$, $R \le
    (T-t)^{-1/2} S$ and~\ref{lem:identities}\ref{i:id:bounds} yields $0 < c < \infty$ with
    \begin{equation}
        \label{eq:reg:ratio}
        \|\uD^j F(t,x)\| \le c (1 + |x|)^{2j} S(t,x) \,, \qquad
        \|\uD^j \rho(t,x)\| \le c (1 + |x|)^{2j} \rho(t,x)
    \end{equation}
    for $F \in \{ U, W, R, \Psi, S \}$, $t \in J$, $x \in \R^{\adim}$, and $j \le j_0$, for any
    given~$j_0$. By the chain rule~\cite[3.1.11]{Federer1969}, $\uD^j \log S$ is a~polynomial in
    the maps $S^{-1} \uD^i S$ with $1 \le i \le j$, and likewise for~$\log \rho$;
    by~\eqref{eq:reg:ratio}, they have polynomial growth. By~\ref{lem:heat_extension}
    \ref{i:he:gauss}, $|\log S(t,x)| + |\log \rho(t,x)| \le c (1 + |x|)^2$ for $t \in J$. Hence
    $\log S$, $\log \rho$ and $\log q = \log \rho - \log S$ have polynomial growth before~$T$; so
    does $q = \exp(\log q) \le 1$, whose derivatives are $q$ times polynomials in the derivatives
    of $\log q$. Similarly the derivatives of $S^{-1}$ are $S^{-1}$ times polynomials in the
    derivatives of $\log S$. By the Leibniz formula and~\eqref{eq:reg:ratio}, $B = S^{-1} W$,
    $\eta = S^{-1} \Psi$, $\kappa = 2^{-1} (T-t)^{-1/2} S^{-1} R$ and $\zeta = \grad \log S$ have
    polynomial growth before~$T$. The final statement follows, because the hypotheses
    of~\ref{lem:ibp} on~$\rho$ concern only a~fixed time.
\end{proof}

\begin{lemma}
    \label{lem:heat_calculus}
    Assume (S). Suppose $Z$ is a~finite dimensional inner product space, $f : \Omega_T \to Z$ and
    $g, \varphi : \Omega_T \to \R$ are of class~$2$, $\varphi > 0$, and $\Phi : (0,\infty) \to \R$
    is of class~$2$. Then the following five identities hold on~$\Omega_T$.
    \begin{enumerate}
        \item \label{i:hc:S} $\heatop(S f) = S \heatop_\zeta f - \kappa S f$.
        \item \label{i:hc:product} $\heatop_\zeta(g f) = g \heatop_\zeta f + (\heatop_\zeta g) f + 2
        \uD f_t(x)(\grad g_t(x))$.
        \item \label{i:hc:square} $\heatop_\zeta |f|^2 = 2 f \bullet \heatop_\zeta f + 2 |\uD f_t
        (x)|^2$.
        \item \label{i:hc:chain} $\heatop_\zeta(\Phi \circ \varphi) = (\Phi' \circ \varphi)
        \heatop_\zeta \varphi + (\Phi'' \circ \varphi) |\grad \varphi_t(x)|^2$.
        \item \label{i:hc:Bq} $\heatop_\zeta B = \kappa B$ and $\heatop_\zeta q = \kappa q$.
    \end{enumerate}
\end{lemma}

\begin{proof}
    The operators $\partial_t$ and $[f \mapsto \uD f_t(x)(\zeta(t,x))]$ are derivations, while
    $\Delta(g f) = g \Delta f + f \Delta g + 2 \uD f(\grad g)$, $\Delta |f|^2 = 2 f \bullet \Delta f
    + 2 |\uD f|^2$, and $\Delta (\Phi \circ \varphi) = (\Phi' \circ \varphi) \Delta \varphi +
    (\Phi'' \circ \varphi) |\grad \varphi|^2$ by~\cite[3.1.11]{Federer1969}. This
    gives~\ref{i:hc:product}--\ref{i:hc:chain}, and also $\heatop(S f) = (\heatop S) f + S \heatop
    f + 2 \uD f(\grad S)$; with~\ref{lem:identities}\ref{i:id:heat} and $\grad S = S \zeta$ the
    latter gives~\ref{i:hc:S}. Since $W = S B$, $\rho = S q$ and $\heatop W = 0 = \heatop
    \rho$, \ref{i:hc:S} yields $S \heatop_\zeta B - \kappa S B = 0$ and $S \heatop_\zeta q - \kappa
    S q = 0$, which is~\ref{i:hc:Bq}.
\end{proof}

\begin{miniremark}
    \label{mr:strategy31}
    Lemma~\ref{lem:DGPR31} is~\cite[Lemma~3.1]{DePhilippisGennaioliPigatiRindler2026}, which
    replaces~\cite[Lemma~3.2]{Gennaioli2026}. The function $J(t) = \textint{}{} \rho_t \log q_t
    \ud \LM^{\adim}$ is bounded between $\gamma \log(\gamma / (1 + (T-t)^{1/2} \delta))$ and~$0$ by
    Jensen's inequality, and its derivative dominates $\textint{}{} \rho_t |\grad \log q_t|^2 \ud
    \LM^{\adim}$. This controls the gradient of~$B$ through the time derivative of $\textint{}{}
    \rho_t |B_t|^2 \ud \LM^{\adim}$, and then $B(\zeta)$ through
    \ref{lem:identities}\ref{i:id:divB}, which is~\cite[(3.9)]{Gennaioli2026}.
\end{miniremark}

\begin{lemma}
    \label{lem:DGPR31}
    Assume (S). There is $0 < C < \infty$, depending only on $\adim$ and~$\Lambda$, such that
    \begin{displaymath}
        \int_0^T \textint{}{} \rho_t \bigl( |\uD B_t|^2 + |B_t(\grad \log \rho_t)|^2 \bigr) \ud
        \LM^{\adim} \ud \LM^1 t \le C \gamma \Bigl( 1 + \log \frac{1 + T^{1/2} \delta}{\gamma}
        \Bigr) \,.
    \end{displaymath}
\end{lemma}

\begin{proof}
    Set $L = \log((1 + T^{1/2} \delta)/\gamma)$; then $L \ge 0$, since $\gamma \le 1$.
    By~\ref{lem:regularity} and~\ref{lem:classes}\ref{i:cl:product}, the maps $\rho \log q$,
    $\rho |B|^2$ and $\rho \, |\uD B_t(x)|^2$ have Gaussian decay before~$T$. For $0 \le t < T$
    define $J(t) = \textint{}{} \rho_t \log q_t \ud \LM^{\adim}$. Since $q \le 1$, we have $J(t)
    \le 0$. Applying Jensen's inequality~\cite[2.4.19]{Federer1969} to the probability measure
    $(1 + (T-t)^{1/2} \delta)^{-1} \LM^{\adim} \restrict S_t$ (see~\ref{lem:identities}
    \ref{i:id:mass}), the convex function $[r \mapsto r \log r]$ on $[0,\infty)$, and the
    function~$q_t$, we obtain
    \begin{equation}
        \label{eq:31:J}
        \gamma \log \frac{\gamma}{1 + (T-t)^{1/2} \delta} \le J(t) \le 0
        \quad \text{for $0 \le t < T$} \,.
    \end{equation}
    By~\ref{lem:heat_calculus}\ref{i:hc:S}\ref{i:hc:chain}\ref{i:hc:Bq} with $f = q \log q$ and
    $\Phi(r) = r \log r$,
    \begin{displaymath}
        \heatop(\rho \log q) = S \bigl( (1 + \log q) \kappa q + q^{-1} |\grad q|^2 \bigr) - \kappa
        S q \log q = \kappa \rho + \rho |\grad \log q|^2 \,.
    \end{displaymath}
    By~\ref{lem:classes}\ref{i:cl:dt}\ref{i:cl:div}, $J'(t) = \textint{}{} \heatop(\rho \log
    q)(t,x) \ud \LM^{\adim} x$. Since $\kappa \ge 0$, integration over $[0,\tau]$,
    \eqref{eq:31:J}, and Lebesgue's increasing convergence theorem~\cite[2.4.7]{Federer1969} as
    $\tau \to T-$ yield
    \begin{equation}
        \label{eq:31:gradq}
        \int_0^T \textint{}{} \rho_t |\grad \log q_t|^2 \ud \LM^{\adim} \ud \LM^1 t \le - J(0)
        \le \gamma L \,.
    \end{equation}
    By~\ref{lem:heat_calculus}\ref{i:hc:S}\ref{i:hc:product}\ref{i:hc:square}\ref{i:hc:Bq},
    $\heatop_\zeta(q |B|^2) = 3 \kappa q |B|^2 + 2 q |\uD B|^2 + 2 \grad q \bullet \grad |B|^2$,
    and since $S \grad q = \rho \grad \log q$,
    \begin{displaymath}
        \heatop(\rho |B|^2) = 2 \rho |\uD B|^2 + 2 \rho \grad \log q \bullet \grad |B|^2 + 2 \kappa
        \rho |B|^2 \,.
    \end{displaymath}
    Since $\grad |B_t|^2(x) \bullet v = 2 B(t,x) \bullet \uD B_t(x)(v)$, we have $|\grad |B_t|^2(x)|
    \le 2 \Lambda |\uD B_t(x)|$ by~\ref{lem:identities}\ref{i:id:bounds}. Hence,
    by~\ref{lem:classes}\ref{i:cl:dt}\ref{i:cl:div} and the inequality $4 \Lambda a b \le a^2 + 4
    \Lambda^2 b^2$,
    \begin{displaymath}
        \frac{\ud}{\ud t} \textint{}{} \rho_t |B_t|^2 \ud \LM^{\adim} \ge \textint{}{} \rho_t |\uD
        B_t|^2 \ud \LM^{\adim} - 4 \Lambda^2 \textint{}{} \rho_t |\grad \log q_t|^2 \ud \LM^{\adim}
        \,.
    \end{displaymath}
    Integrating over $[0,\tau]$, using $\textint{}{} \rho_\tau |B_\tau|^2 \ud \LM^{\adim} \le
    \Lambda^2 \gamma$ and~\eqref{eq:31:gradq}, and letting $\tau \to T-$, we obtain
    \begin{equation}
        \label{eq:31:DB}
        \int_0^T \textint{}{} \rho_t |\uD B_t|^2 \ud \LM^{\adim} \ud \LM^1 t \le \Lambda^2 \gamma +
        4 \Lambda^2 \gamma L \,.
    \end{equation}
    Since $\rho \|\eta\|^2 \le (T-t)^{-1} \rho$, and since $\rho \|\eta\|^2 \le S \|\eta\|^2 =
    S^{-1} \|\Psi\|^2 \le (S^{-1} R) R \le (T-t)^{-1/2} R$ by~\ref{lem:identities}
    \ref{i:id:div}\ref{i:id:bounds}, statement~\ref{lem:identities}\ref{i:id:mass} gives
    \begin{equation}
        \label{eq:31:eta}
        \textint{}{} \rho_t \|\eta_t\|^2 \ud \LM^{\adim} \le m(t) \,, \quad \text{where} \quad m(t)
        = \min \{ \gamma (T-t)^{-1}, \delta (T-t)^{-1/2} \} \,, \quad \text{for $0 \le t < T$} \,.
    \end{equation}
    We show
    \begin{equation}
        \label{eq:31:min}
        \int_0^T m(t) \ud \LM^1 t \le 4 \gamma \log(1 + T^{1/2} \delta / \gamma) \le 4 \gamma L
        \,.
    \end{equation}
    Set $a = T^{1/2} \delta / \gamma$. If $a \le 1$, then $\int_0^T m \ud \LM^1 \le \int_0^T
    \delta (T-t)^{-1/2} \ud \LM^1 t = 2 \gamma a \le 4 \gamma \log(1 + a)$, since $\log(1 + a) \ge
    a/2$ for $0 \le a \le 1$. If $a > 1$, then $\gamma^2 \delta^{-2} < T$ and, with $s = T - t$,
    \begin{displaymath}
        \int_0^T m \ud \LM^1 \le \int_0^{\gamma^2/\delta^2} \delta s^{-1/2} \ud \LM^1 s +
        \int_{\gamma^2/\delta^2}^T \gamma s^{-1} \ud \LM^1 s = 2 \gamma + 2 \gamma \log a \le 4
        \gamma \log(1 + a) \,,
    \end{displaymath}
    since $[a \mapsto 4 \log(1+a) - 2 - 2 \log a]$ is nondecreasing on $[1,\infty)$ and positive
    at~$1$. The second inequality of~\eqref{eq:31:min} holds because $1 + a \le (1 + T^{1/2}
    \delta)/\gamma$.

    By~\ref{lem:identities}\ref{i:id:divB}, $|B(\zeta)| = \sup \{ B(\zeta) \bullet y : |y| \le 1 \}
    \le \|\eta\| + \|\mathrm{div}\, B\|$, and $\|\mathrm{div}\, B_t(x)\| \le \adim^{1/2} |\uD
    B_t(x)|$ by the Cauchy--Schwarz inequality. Moreover, $B(\grad \log \rho) = B(\zeta) +
    B(\grad \log q)$ and $\|B\| \le \Lambda$. Hence
    \begin{displaymath}
        |B(\grad \log \rho)|^2 \le 2 |B(\zeta)|^2 + 2 \Lambda^2 |\grad \log q|^2 \le 4 \|\eta\|^2
        + 4 \adim |\uD B|^2 + 2 \Lambda^2 |\grad \log q|^2 \,.
    \end{displaymath}
    Combining this with \eqref{eq:31:gradq}, \eqref{eq:31:DB}, \eqref{eq:31:eta}
    and~\eqref{eq:31:min}, we obtain the conclusion with $C = (16 \adim + 6) \Lambda^2 + 16$.
\end{proof}

\begin{miniremark}
    \label{mr:strategy32}
    Lemma~\ref{lem:DGPR32} is~\cite[Lemma~3.2]{DePhilippisGennaioliPigatiRindler2026}. At each
    time $t$ the field $Q = \hat\chi(\transpose{B_t} \circ B_t)$ is an~approximate projection of
    rank at least~$\vdim$ on the set where $U_t \ge (T-t)^{1/2} R_t$. Along its image the Hessian
    of $\log \rho_t$ is estimated by integration by parts~\ref{lem:ibp} and~\ref{lem:DGPR31}; in
    the remaining $\adim - \vdim$ directions it is estimated by the log-semiconvexity
    \ref{lem:semiconvex}. The trace comparison~\ref{lem:trace} combines the two estimates.
\end{miniremark}

\begin{lemma}
    \label{lem:DGPR32}
    Assume (S). There is $0 < C < \infty$, depending only on $\adim$, $\lambda$ and~$\Lambda$,
    such that
    \begin{displaymath}
        \int_0^\tau \Bigl( I(\rho_t) - \frac{(\adim - \vdim) \gamma}{2 (T-t)} \Bigr) \ud \LM^1 t
        \le C \gamma \Bigl( 1 + \log \frac{1 + T^{1/2} \delta}{\gamma} \Bigr)
        \quad \text{for $0 \le \tau < T$} \,.
    \end{displaymath}
\end{lemma}

\begin{proof}
    Define $\chi : \R \to [0,1]$ by $\chi(a) = 0$ for $a \le \lambda^2/16$, $\chi(a) = 16 a /
    (3 \lambda^2) - 1/3$ for $\lambda^2/16 \le a \le \lambda^2/4$, and $\chi(a) = 1$ for $a \ge
    \lambda^2/4$. Then $\Lip(\chi) = 16/(3 \lambda^2)$ and
    \begin{equation}
        \label{eq:32:chi}
        \chi(a)^2 \le \chi(a) \le 16 a / (3 \lambda^2) \quad \text{for $a \ge 0$} \,.
    \end{equation}
    Fix $0 \le t < T$ and define $Q(x) = \hat\chi(\transpose{B(t,x)} \circ B(t,x))$ for $x \in
    \R^{\adim}$, where $\transpose{B(t,x)} \in \Hom(Y,\R^{\adim})$. Let $x \in \R^{\adim}$, let
    $u_1, \ldots, u_{\adim}$ be an~orthonormal basis with $\transpose{B(t,x)} \circ B(t,x)(u_i) =
    a_i u_i$, and note $a_i = |B(t,x)(u_i)|^2 \ge 0$. By~\ref{def:spectral}, $Q(x) \in
    \SymEnd{\R^{\adim}}$, $\|Q(x)\| \le 1$, $|Q(x)|^2 = \textsum{i}{} \chi(a_i)^2 \le \adim$, and,
    by~\eqref{eq:32:chi},
    \begin{equation}
        \label{eq:32:QB}
        |Q(x)(w)|^2 = \textsum{i}{} \chi(a_i)^2 (u_i \bullet w)^2 \le \frac{16}{3 \lambda^2}
        \textsum{i}{} a_i (u_i \bullet w)^2 = \frac{16}{3 \lambda^2} |B(t,x)(w)|^2
        \quad \text{for $w \in \R^{\adim}$} \,.
    \end{equation}
    If $U(t,x) \ge (T-t)^{1/2} R(t,x)$, then $\card \{ i : a_i \ge \lambda^2/4 \} \ge \vdim$
    by~\ref{lem:identities}\ref{i:id:sv} and~\ref{lem:sv}\ref{i:sv:eigen}, hence $|Q(x)|^2 \ge
    \vdim$.

    By~\ref{lem:spectral}, $|Q(x) - Q(x')| \le \Lip(\chi) |\transpose{B(t,x)} \circ B(t,x) -
    \transpose{B(t,x')} \circ B(t,x')|$. Hence $Q$ is locally Lipschitzian, and at each point~$x$
    of differentiability of~$Q$, using $|\transpose{(\uD B_t(x)(v))} \circ B(t,x) +
    \transpose{B(t,x)} \circ \uD B_t(x)(v)| \le 2 \Lambda |\uD B_t(x)(v)|$ for $v \in \R^{\adim}$,
    \begin{equation}
        \label{eq:32:DQ}
        |\uD Q(x)| \le 2 \Lambda \Lip(\chi) |\uD B_t(x)| = \frac{32 \Lambda}{3 \lambda^2} |\uD
        B_t(x)| \,.
    \end{equation}
    By~\ref{lem:regularity}, the maps $e_j = [x \mapsto Q(x)(\varepsilon_j)]$, for $j = 1, \ldots,
    \adim$, are locally Lipschitzian maps of polynomial growth.

    Let $h(x) = \uD^2 \log \rho_t(x)$. By~\ref{lem:semiconvex}, $h(x)(v,v) \ge -(2(T-t))^{-1}
    |v|^2$. Applying~\ref{lem:trace} to $g = h(x) + (2(T-t))^{-1} [(v,w) \mapsto v \bullet w]$ and
    $Q(x)$, we obtain
    \begin{equation}
        \label{eq:32:split}
        - \Delta \log \rho_t(x) \le \frac{\adim - |Q(x)|^2}{2 (T-t)} - \textsum{j=1}{\adim}
        h(x)(e_j(x),e_j(x)) \quad \text{for $x \in \R^{\adim}$} \,.
    \end{equation}
    Let $G = \R^{\adim} \cap \{ x : U(t,x) \ge (T-t)^{1/2} R(t,x) \}$. Since $\rho_t \le U_t <
    (T-t)^{1/2} R_t$ on $\R^{\adim} \without G$, statement~\ref{lem:identities}\ref{i:id:mass}
    gives $\textint{\R^{\adim} \without G}{} \rho_t \ud \LM^{\adim} \le \min \{ \gamma,
    (T-t)^{1/2} \delta \}$. Since $\adim - |Q|^2 \le \adim - \vdim$ on~$G$ and $\adim - |Q|^2 \le
    \adim$ everywhere,
    \begin{equation}
        \label{eq:32:trace}
        \textint{}{} \rho_t (\adim - |Q|^2) \ud \LM^{\adim} \le (\adim - \vdim) \gamma + \adim \min
        \{ \gamma, (T-t)^{1/2} \delta \} \,.
    \end{equation}
    Since $\grad \rho_t = \rho_t \grad \log \rho_t$ has Gaussian decay, \ref{lem:classes}
    \ref{i:cl:div} applied to $\mathrm{div}(\rho_t \grad \log \rho_t) = \rho_t \Delta \log \rho_t +
    \rho_t |\grad \log \rho_t|^2$ gives $I(\rho_t) = - \textint{}{} \rho_t \Delta \log \rho_t \ud
    \LM^{\adim}$. By~\ref{lem:ibp},
    \begin{displaymath}
        - \textint{}{} \rho_t \, h(e_j,e_j) \ud \LM^{\adim} = \textint{}{} \rho_t \bigl[
        (\mathrm{div}\, e_j + \uD \log \rho_t(e_j))^2 - \trace(\uD e_j \circ \uD e_j) \bigr] \ud
        \LM^{\adim} \,.
    \end{displaymath}
    At points of differentiability of~$Q$ we have $(\mathrm{div}\, e_j)^2 \le \adim |\uD e_j|^2$,
    $|\trace(\uD e_j \circ \uD e_j)| \le |\uD e_j|^2$, $\textsum{j}{} |\uD e_j|^2 = |\uD Q|^2$,
    and, since $Q$ is symmetric, $\textsum{j}{} (\uD \log \rho_t(e_j))^2 = \textsum{j}{} (Q(\grad
    \log \rho_t) \bullet \varepsilon_j)^2 = |Q(\grad \log \rho_t)|^2$. Combining these facts with
    \eqref{eq:32:split}, \eqref{eq:32:trace}, \eqref{eq:32:QB} and~\eqref{eq:32:DQ}, we obtain
    \begin{equation}
        \label{eq:32:I}
        I(\rho_t) - \frac{(\adim - \vdim) \gamma}{2 (T-t)} \le \frac{\adim}{2} m(t) + C_1
        \textint{}{} \rho_t \bigl( |\uD B_t|^2 + |B_t(\grad \log \rho_t)|^2 \bigr) \ud
        \LM^{\adim} \,,
    \end{equation}
    where $m$ is as in~\ref{lem:DGPR31}\eqref{eq:31:eta} and $C_1 = (2 \adim + 1)(32 \Lambda / (3
    \lambda^2))^2 + 32/(3 \lambda^2)$. The right-hand side of~\eqref{eq:32:I} and $I(\rho_t)$ are
    continuous in~$t$ by~\ref{lem:regularity} and~\ref{lem:classes}\ref{i:cl:dt}. Integrating
    \eqref{eq:32:I} over $[0,\tau]$ and applying~\ref{lem:DGPR31}\eqref{eq:31:min}
    and~\ref{lem:DGPR31} gives the conclusion.
\end{proof}

\section{Entropy}
\label{sec:entropy}

\begin{miniremark}
    \label{mr:sec_entropy}
    This section reproduces the two entropy bounds~\cite[(4.1), (4.2)]
    {DePhilippisGennaioliPigatiRindler2026}. The upper bound~\ref{lem:upper_entropy} follows
    from~\ref{lem:DGPR32}, since the time derivative of the entropy is the Fisher
    information~\ref{lem:entropy}. The lower bound~\ref{lem:lower_entropy} holds when $\mu_0 = \mu
    \restrict E$ for a~compact $\vdim$ rectifiable set~$E$; it follows from the Gaussian tube
    estimate~\ref{lem:tube}, which rests on the Minkowski content
    theorem~\cite[3.2.39]{Federer1969}. The hypotheses (S) of~\ref{mr:standing} remain in force.
\end{miniremark}

\begin{definition}
    \label{def:entropy}
    Assume (S). For $t < T$ we define $H(\rho_t) = \textint{}{} \rho_t \log \rho_t \ud
    \LM^{\adim}$ and
    \begin{displaymath}
        \tilde H(\rho_t) = H(\rho_t) + 2^{-1} (\adim - \vdim) \gamma \log(T-t) \,.
    \end{displaymath}
\end{definition}

\begin{lemma}
    \label{lem:entropy}
    Assume (S). Then $[t \mapsto H(\rho_t)]$ is of class~$1$ on $\R \cap \{ t : t < T \}$ with
    derivative $[t \mapsto I(\rho_t)]$, and $H(\rho_0) \le \gamma \log \gamma - 2^{-1} \adim
    \gamma \log(4 \pi T)$.
\end{lemma}

\begin{proof}
    By~\ref{lem:regularity} and~\ref{lem:classes}\ref{i:cl:product}, the map $\rho \log \rho$ and
    the map $[(t,x) \mapsto (1 + \log \rho(t,x)) \grad \rho_t(x)]$ have Gaussian decay
    before~$T$. Since $\partial_t \rho = - \Delta \rho$ by~\ref{lem:identities}\ref{i:id:heat},
    \begin{displaymath}
        \partial_t (\rho \log \rho) = - (1 + \log \rho) \Delta \rho = - \mathrm{div}\bigl[ (1 +
        \log \rho) \grad \rho \bigr] + \rho |\grad \log \rho|^2 \,.
    \end{displaymath}
    Hence the first assertion follows from~\ref{lem:classes}\ref{i:cl:dt}\ref{i:cl:div}.
    By~\ref{lem:kernel}\ref{i:k:max}, $\rho_0 \le \gamma (4 \pi T)^{-\adim/2}$, hence $\rho_0 \log
    \rho_0 \le \rho_0 \log(\gamma (4 \pi T)^{-\adim/2})$; integration and~\ref{lem:identities}
    \ref{i:id:mass} give the second assertion.
\end{proof}

\begin{lemma}
    \label{lem:upper_entropy}
    Assume (S). There is $0 < C < \infty$, depending only on $\adim$, $\lambda$ and~$\Lambda$,
    such that
    \begin{displaymath}
        \tilde H(\rho_\tau) \le - 2^{-1} \vdim \gamma \log T + C \gamma \Bigl( 1 + \log \frac{1 +
        T^{1/2} \delta}{\gamma} \Bigr)
        \quad \text{for $0 \le \tau < T$} \,.
    \end{displaymath}
\end{lemma}

\begin{proof}
    By~\ref{lem:entropy}, $\tilde H(\rho_\tau) - \tilde H(\rho_0) = \int_0^\tau ( I(\rho_t) -
    (\adim - \vdim) \gamma / (2(T-t)) ) \ud \LM^1 t$, which is bounded by~\ref{lem:DGPR32}.
    By~\ref{lem:entropy}, $\gamma \le 1$ and $4 \pi > 1$,
    \begin{displaymath}
        \tilde H(\rho_0) \le \gamma \log \gamma - 2^{-1} \adim \gamma \log(4 \pi) - 2^{-1} \adim
        \gamma \log T + 2^{-1} (\adim - \vdim) \gamma \log T \le - 2^{-1} \vdim \gamma \log T \,.
        \qedhere
    \end{displaymath}
\end{proof}

\begin{lemma}
    \label{lem:tube}
    Suppose $E \subseteq \R^{\adim}$ is compact and $\vdim$~rectifiable, $\HM^{\vdim}(E) > 0$,
    and
    \begin{displaymath}
        \gamma_s = \textint{}{} \exp \bigl( - \dist(x,E)^2 / (4s) \bigr) \ud \LM^{\adim} x
        \quad \text{for $0 < s < \infty$} \,.
    \end{displaymath}
    Then there is $0 < s_0 \le 1$ such that $\gamma_s \le c(\adim) \, s^{(\adim - \vdim)/2}
    \HM^{\vdim}(E)$ for $0 < s < s_0$, where $c(\adim) = 1 + 2 (4 \pi)^{\adim/2}$.
\end{lemma}

\begin{proof}
    Let $E_r = \R^{\adim} \cap \{ x : \dist(x,E) < r \}$ for $0 < r < \infty$. Since
    \begin{displaymath}
        \exp(-d^2 / (4s)) = \int_0^\infty \CF{(d,\infty)}(r) \, (2s)^{-1} r \exp(-r^2/(4s)) \ud
        \LM^1 r \quad \text{for $0 \le d < \infty$} \,,
    \end{displaymath}
    Fubini's theorem~\cite[2.6.2]{Federer1969} gives
    \begin{equation}
        \label{eq:tube:layer}
        \gamma_s = \int_0^\infty (2s)^{-1} r \exp(-r^2/(4s)) \, \LM^{\adim}(E_r) \ud \LM^1 r \,.
    \end{equation}
    Since $E$ is closed and $\vdim$~rectifiable, $\lim_{r \to 0+} \LM^{\adim}(E_r) /
    (\unitmeasure{\adim - \vdim} r^{\adim - \vdim}) = \HM^{\vdim}(E)$ by~\cite[3.2.37,
    3.2.39]{Federer1969}. Since $\HM^{\vdim}(E) > 0$, there is $0 < r_1 \le 1$ with
    $\LM^{\adim}(E_r) \le 2 \unitmeasure{\adim - \vdim} r^{\adim - \vdim} \HM^{\vdim}(E)$ for $0
    < r < r_1$. Choose $0 < R_0 < \infty$ with $E \subseteq \cball{0}{R_0}$; then
    $\LM^{\adim}(E_r) \le \unitmeasure{\adim} (R_0 + r)^{\adim}$. With $m = \adim - \vdim$, the
    substitution $r = 2 s^{1/2} \tau$ and $\unitmeasure{m} \Gamma(1 + m/2) = \pi^{m/2}$ give
    \begin{align*}
        \int_0^{r_1} (2s)^{-1} r \exp(-r^2/(4s)) \, 2 \unitmeasure{m} r^m \HM^{\vdim}(E) \ud \LM^1
        r & \le 2 \unitmeasure{m} \, 2^m \Gamma(1 + m/2) \, s^{m/2} \HM^{\vdim}(E) \\
        & \le 2 (4 \pi)^{\adim/2} s^{m/2} \HM^{\vdim}(E) \,.
    \end{align*}
    For $0 < s \le 1$, the substitution $r = s^{1/2} \tau$ gives
    \begin{displaymath}
        \int_{r_1}^\infty (2s)^{-1} r \exp(-r^2/(4s)) \, \unitmeasure{\adim} (R_0 + r)^{\adim} \ud
        \LM^1 r \le c_1 \exp(-r_1^2/(8s)) \,,
    \end{displaymath}
    where $c_1 = \unitmeasure{\adim} \int_0^\infty 2^{-1} \tau \exp(-\tau^2/8) (R_0 +
    \tau)^{\adim} \ud \LM^1 \tau$.
    Since $s^{-m/2} \exp(-r_1^2/(8s)) \to 0$ as $s \to 0+$, there is $0 < s_0 \le 1$ with $c_1
    \exp(-r_1^2/(8s)) \le s^{m/2} \HM^{\vdim}(E)$ for $0 < s < s_0$. The conclusion follows
    from~\eqref{eq:tube:layer} and the last two displays.
\end{proof}

\begin{lemma}
    \label{lem:lower_entropy}
    Assume (S). Suppose $E \subseteq \R^{\adim}$ is compact and $\vdim$~rectifiable,
    $\HM^{\vdim}(E) > 0$, and $\mu_0 = \mu \restrict E$. Then
    \begin{displaymath}
        \liminf_{t \to T-} \tilde H(\rho_t) \ge \gamma \log \frac{\gamma}{\HM^{\vdim}(E)} - \gamma
        \bigl( \log c(\adim) + 2^{-1} \adim \bigr) \,,
    \end{displaymath}
    where $c(\adim)$ is as in~\ref{lem:tube}.
\end{lemma}

\begin{proof}
    Let $s_0$ and $\gamma_s$ be as in~\ref{lem:tube}, and suppose $T - s_0 < t < T$ and $t \ge 0$.
    Set $s = T - t$, $\varphi = \gamma^{-1} \rho_t$ and $\psi(x) = \gamma_s^{-1} \exp(-\dist(x,E)^2
    / (4s))$ for $x \in \R^{\adim}$. Both functions are positive and continuous, and their
    integrals equal~$1$ by~\ref{lem:identities}\ref{i:id:mass}. Jensen's
    inequality~\cite[2.4.19]{Federer1969}, applied to the probability measure $\LM^{\adim}
    \restrict \psi$, the convex function $[r \mapsto r \log r]$ and the function $\varphi/\psi$,
    gives $\textint{}{} \varphi \log(\varphi/\psi) \ud \LM^{\adim} \ge 0$; the integrals exist
    by~\ref{lem:regularity}, since $\dist(x,E)^2 \le 2 |x|^2 + 2 R_0^2$ for some $R_0$. Hence
    \begin{displaymath}
        H(\rho_t) = \gamma \textint{}{} \varphi \log \varphi \ud \LM^{\adim} + \gamma \log \gamma
        \ge \gamma \log \gamma - \gamma \log \gamma_s - (4s)^{-1} \textint{}{} \rho_t(x)
        \dist(x,E)^2 \ud \LM^{\adim} x \,.
    \end{displaymath}
    Since $\mu_0(\R^{\adim} \without E) = 0$ and $\dist(x,E) \le |x - y|$ for $y \in E$, Fubini's
    theorem~\cite[2.6.2]{Federer1969} and~\ref{lem:kernel}\ref{i:k:moment} give
    \begin{displaymath}
        \textint{}{} \rho_t(x) \dist(x,E)^2 \ud \LM^{\adim} x \le \textint{E}{} \textint{}{} |x -
        y|^2 p_s(x-y) \ud \LM^{\adim} x \ud \mu_0 y = 2 \adim s \gamma \,.
    \end{displaymath}
    By~\ref{lem:tube}, $\log \gamma_s \le \log c(\adim) + 2^{-1}(\adim - \vdim) \log s + \log
    \HM^{\vdim}(E)$. Therefore
    \begin{displaymath}
        \tilde H(\rho_t) = H(\rho_t) + 2^{-1} (\adim - \vdim) \gamma \log s \ge \gamma \log
        \frac{\gamma}{\HM^{\vdim}(E)} - \gamma \log c(\adim) - 2^{-1} \adim \gamma \,.
        \qedhere
    \end{displaymath}
\end{proof}

\section{The mass bound}
\label{sec:main}

\begin{miniremark}
    \label{mr:sec_main}
    This section proves~\cite[Theorem~1.1]{DePhilippisGennaioliPigatiRindler2026} for countably
    $(\HM^{\vdim},\vdim)$ rectifiable sets~$M$. A~compact $\vdim$ rectifiable set $E \subseteq M$
    carrying at least half of~$\mu$ is provided by~\ref{lem:inner}. For compactly supported~$\mu$,
    the entropy bounds~\ref{lem:upper_entropy} and~\ref{lem:lower_entropy} with $T =
    \|\Upsilon\|(\R^{\adim})^{-2}$ give the conclusion~\ref{lem:compact_case}. The general case
    follows by a~cut-off in~\ref{thm:main}. The hypotheses (S) of~\ref{mr:standing} are not
    assumed in this section.
\end{miniremark}

\begin{lemma}
    \label{lem:inner}
    Suppose $M \subseteq \R^{\adim}$ is a~Borel set which is countably $(\HM^{\vdim},\vdim)$
    rectifiable, $\HM^{\vdim}(M) < \infty$, $\theta : M \to [0,\infty)$ is $\HM^{\vdim} \restrict
    M$ summable, $\mu = (\HM^{\vdim} \restrict M) \restrict \theta$, and $0 < \epsilon < \infty$.
    Then there is a~compact $\vdim$~rectifiable set $E \subseteq M$ with $\mu(\R^{\adim} \without
    E) < \epsilon$.
\end{lemma}

\begin{proof}
    By~\cite[3.2.14]{Federer1969} there are bounded sets $A_1, A_2, \ldots \subseteq \R^{\vdim}$
    and Lipschitzian maps $f_i : A_i \to \R^{\adim}$ with $\HM^{\vdim}(M \without
    \tbcup_{i=1}^\infty f_i[A_i]) = 0$. Extend $f_i$ to a~Lipschitzian map of~$\R^{\vdim}$ into
    $\R^{\adim}$ by Kirszbraun's theorem~\cite[2.10.43]{Federer1969}, and set $F_i = f_i[\Clos
    A_i]$. Then $F_i$ is compact and $\vdim$~rectifiable, and $\mu(\R^{\adim} \without
    \tbcup_{i=1}^\infty F_i) = 0$. Since $\mu(\R^{\adim}) < \infty$, there is $j$ with
    $\mu(\R^{\adim} \without F) < \epsilon/2$, where $F = \tbcup_{i=1}^j F_i$ is compact and
    $\vdim$~rectifiable~\cite[3.2.14]{Federer1969}. The set $M \cap F$ is a~Borel set and $\mu$ is
    a~Borel measure; by~\cite[2.2.2]{Federer1969} there is a~closed set $E \subseteq M \cap F$
    with $\mu((M \cap F) \without E) < \epsilon/2$. The set~$E$ is compact, and it is
    $\vdim$~rectifiable as a~closed subset of~$F$~\cite[3.2.14]{Federer1969}. Since $\mu(\R^{\adim}
    \without M) = 0$, we have $\mu(\R^{\adim} \without E) < \epsilon$.
\end{proof}

\begin{miniremark}
    \label{mr:hypotheses_main}
    In~\ref{lem:compact_case} and~\ref{thm:main} we consider the following data. The space $Y$
    is a~finite dimensional inner product space, $K \subseteq \Hom(\R^{\adim},Y)$ is nonempty,
    compact and convex, $0 < \lambda < \infty$, $0 < \Lambda < \infty$, and $\sv{\vdim}(X) \ge
    \lambda$ and $|X| \le \Lambda$ for $X \in K$. The set $M \subseteq \R^{\adim}$ is a~Borel set
    which is countably $(\HM^{\vdim},\vdim)$ rectifiable, $\HM^{\vdim}(M) < \infty$, $\theta : M
    \to [0,\infty)$ is $\HM^{\vdim} \restrict M$ summable, $X : M \to K$ is $\HM^{\vdim}
    \restrict M$ measurable, $\mu = (\HM^{\vdim} \restrict M) \restrict \theta$, and
    \begin{displaymath}
        \sigma(\Phi) = \textint{}{} X(x) \bullet \Phi(x) \ud \mu x
        \quad \text{for $\Phi \in \ccspace{\R^{\adim},\Hom(\R^{\adim},Y)}$} \,.
    \end{displaymath}
    The distribution $\Upsilon \in \distr{\R^{\adim}}{Y}$ is representable by integration,
    $\|\Upsilon\|(\R^{\adim}) < \infty$, and $\Upsilon(g) = - \sigma(\uD g)$ for $g \in
    \dspace{\R^{\adim}}{Y}$.
\end{miniremark}

\begin{lemma}
    \label{lem:compact_case}
    There is $0 < C < \infty$, depending only on $\adim$, $\lambda$ and~$\Lambda$, with the
    following property. If the data of~\ref{mr:hypotheses_main} are given and $\mu(\R^{\adim}
    \without \cball{0}{r_0}) = 0$ for some $0 < r_0 < \infty$, then
    \begin{displaymath}
        \mu(\R^{\adim}) \le C \, \HM^{\vdim}(M)^{1/\vdim} \, \|\Upsilon\|(\R^{\adim}) \,.
    \end{displaymath}
\end{lemma}

\begin{proof}
    Let $C_0$ be the constant of~\ref{lem:upper_entropy}, $c(\adim)$ that of~\ref{lem:tube}, $C_1 =
    C_0 + \log c(\adim) + 2^{-1} \adim$, $C_2 = \log 2 + C_1 (1 + \log 4)$, and $C = \exp(C_2)$. If
    $\mu(\R^{\adim}) = 0$ there is nothing to prove. Otherwise, replacing $\theta$ by $\theta /
    \mu(\R^{\adim})$ and $\Upsilon$ by $\Upsilon / \mu(\R^{\adim})$, we may assume
    $\mu(\R^{\adim}) = 1$, since both sides of the conclusion are homogeneous of degree one in
    $(\theta,\Upsilon)$. Extend $X$ to~$\R^{\adim}$ by a~constant value in~$K$ on $\R^{\adim}
    \without M$, and set $\delta = \|\Upsilon\|(\R^{\adim})$.

    By~\ref{lem:inner} there is a~compact $\vdim$~rectifiable set $E \subseteq M$ with $\gamma =
    \mu(E) > 1/2$. Since $\mu$ is absolutely continuous with respect to $\HM^{\vdim}$, we have $0
    < \HM^{\vdim}(E) \le \HM^{\vdim}(M) < \infty$. For each $0 < T < \infty$, the hypotheses (S)
    of~\ref{mr:standing} hold with $\mu_0 = \mu \restrict E$. Combining~\ref{lem:upper_entropy}
    with~\ref{lem:lower_entropy} and dividing by~$\gamma$, we obtain
    \begin{equation}
        \label{eq:cc:main}
        \log \frac{\gamma}{\HM^{\vdim}(E)} \le - \frac{\vdim}{2} \log T + C_1 \Bigl( 1 + \log
        \frac{1 + T^{1/2} \delta}{\gamma} \Bigr)
        \quad \text{for $0 < T < \infty$} \,.
    \end{equation}
    If $\delta = 0$, the right-hand side of~\eqref{eq:cc:main} tends to $-\infty$ as $T \to \infty$,
    which is impossible. Hence $\delta > 0$. Taking $T = \delta^{-2}$ in~\eqref{eq:cc:main} and
    using $1/2 < \gamma \le 1$, we obtain $- \log 2 - \log \HM^{\vdim}(E) \le \vdim \log \delta +
    C_1 (1 + \log 4)$, that is, $\HM^{\vdim}(E) \, \delta^{\vdim} \ge \exp(-C_2)$. Hence $1 \le
    \exp(C_2/\vdim) \HM^{\vdim}(E)^{1/\vdim} \delta \le C \, \HM^{\vdim}(M)^{1/\vdim} \delta$.
\end{proof}

\begin{theorem}
    \label{thm:main}
    There is $0 < C < \infty$, depending only on $\adim$, $\lambda$ and~$\Lambda$, such that
    \begin{displaymath}
        \mu(\R^{\adim}) \le C \, \HM^{\vdim}(M)^{1/\vdim} \, \|\Upsilon\|(\R^{\adim})
    \end{displaymath}
    whenever the data of~\ref{mr:hypotheses_main} are given.
\end{theorem}

\begin{proof}
    Let $C$ be the constant of~\ref{lem:compact_case}. Choose $\varphi \in \dspace{\R^{\adim}}{\R}$
    with $0 \le \varphi \le 1$, $\varphi(x) = 1$ for $|x| \le 1$, and $\spt \varphi \subseteq
    \cball{0}{2}$, and set $c_0 = \sup \{ |\grad \varphi(x)| : x \in \R^{\adim} \}$. For $1 \le
    \varrho < \infty$ let $\varphi_\varrho(x) = \varphi(\varrho^{-1} x)$, $\theta_\varrho =
    \varphi_\varrho \theta$, $\mu_\varrho = (\HM^{\vdim} \restrict M) \restrict \theta_\varrho$,
    $\sigma_\varrho(\Phi) = \textint{}{} X \bullet \Phi \ud \mu_\varrho$, and $\Upsilon_\varrho(g)
    = - \sigma_\varrho(\uD g)$ for $g \in \dspace{\R^{\adim}}{Y}$. Since $\uD(\varphi_\varrho g) =
    \varphi_\varrho \uD g + \uD \varphi_\varrho \cdot g$ and $X \bullet (\uD \varphi_\varrho \cdot
    g) = \uD \varphi_\varrho(\transpose{X}(g))$,
    \begin{displaymath}
        \Upsilon_\varrho(g) = \Upsilon(\varphi_\varrho g) + \textint{}{} \uD
        \varphi_\varrho(x)(\transpose{X(x)}(g(x))) \ud \mu x \,.
    \end{displaymath}
    Since $\|\transpose{X(x)}\| \le |X(x)| \le \Lambda$ and $|\grad \varphi_\varrho| \le c_0 /
    \varrho$, we obtain $|\Upsilon_\varrho(g)| \le \textint{}{} |g| \ud \nu_\varrho$ with
    $\nu_\varrho = \|\Upsilon\| + c_0 \Lambda \varrho^{-1} \mu$. By~\ref{mr:functionals},
    $\Upsilon_\varrho$ is representable by integration and $\|\Upsilon_\varrho\|(\R^{\adim}) \le
    \|\Upsilon\|(\R^{\adim}) + c_0 \Lambda \varrho^{-1} \mu(\R^{\adim})$. The data
    $(\theta_\varrho, \mu_\varrho, \sigma_\varrho, \Upsilon_\varrho)$ satisfy the hypotheses
    of~\ref{lem:compact_case} with $r_0 = 2 \varrho$. Hence
    \begin{displaymath}
        \mu_\varrho(\R^{\adim}) \le C \, \HM^{\vdim}(M)^{1/\vdim} \bigl( \|\Upsilon\|(\R^{\adim}) +
        c_0 \Lambda \varrho^{-1} \mu(\R^{\adim}) \bigr) \,.
    \end{displaymath}
    Since $\mu(\cball{0}{\varrho}) \le \mu_\varrho(\R^{\adim})$ and $\mu(\cball{0}{\varrho}) \to
    \mu(\R^{\adim})$ as $\varrho \to \infty$, letting $\varrho \to \infty$ gives the
    conclusion.
\end{proof}

\section{Varifolds and the Michael--Simon inequality}
\label{sec:varifolds}

\begin{miniremark}
    \label{mr:sec_varifolds}
    This section deduces~\cite[Theorem~1.6]{DePhilippisGennaioliPigatiRindler2026}
    from~\ref{thm:main}. It recalls Allard's notation, defines the anisotropic first variation
    and the map~$B_F$, proves the first variation formula~\ref{lem:first_variation} following
    \cite[Lemma~A.2]{DePhilippis2018}, and shows in~\ref{lem:KF} that (AC1) yields the
    constants $\lambda$ and~$\Lambda$ of~\ref{mr:hypotheses_main}.
\end{miniremark}

\begin{miniremark}
    \label{mr:varifolds}
    Following~\cite[2.3]{Allard1972}, we identify each $\vdim$~dimensional linear subspace
    of~$\R^{\adim}$ with the orthogonal projection of~$\R^{\adim}$ onto it, and write
    $\pgrass{\adim}{\vdim}$ for the set of these projections; for $S \in \pgrass{\adim}{\vdim}$ we
    set $S^\perp = \id{\R^{\adim}} - S$. The set $\pgrass{\adim}{\vdim}$ is a~compact submanifold
    of class~$\infty$ of~$\End{\R^{\adim}}$, and
    \begin{displaymath}
        \Tan(\pgrass{\adim}{\vdim},S) = \SymEnd{\R^{\adim}} \cap \{ A : S \circ A \circ S = 0,\
        S^\perp \circ A \circ S^\perp = 0 \}
    \end{displaymath}
    by~\cite[Appendix~A]{DePhilippis2018}. Following~\cite[3.1]{Allard1972}, a~$\vdim$
    dimensional varifold in~$\R^{\adim}$ is a~Radon measure~$V$ over $\R^{\adim} \times
    \pgrass{\adim}{\vdim}$, $\Var{\vdim}(\R^{\adim})$ is the space of these varifolds, and
    $\|V\|$ is the weight of~$V$. The space $\RVar{\vdim}(\R^{\adim})$ of rectifiable varifolds is
    that of~\cite[3.5]{Allard1972}. For a~smooth map $f : \R^{\adim} \to \R^{\adim}$, $f_\# V$ is
    the image varifold of~\cite[3.2(1)]{Allard1972}. We denote by $\density^{\vdim}(\|V\|,x)$,
    $\density^{*\vdim}(\|V\|,x)$ and $\Tan^{\vdim}(\|V\|,x)$ the density, the upper density and
    the approximate tangent plane, as in~\cite[2.10.19]{Federer1969} and~\cite[3.5]{Allard1972}.

    In this section $F : \pgrass{\adim}{\vdim} \to (0,\infty)$ is of class~$1$
    in the sense of~\cite[3.1.22]{Federer1969}, and $\uD F(S) \in
    \Hom(\Tan(\pgrass{\adim}{\vdim},S),\R)$ denotes its differential at~$S$. For $V \in
    \Var{\vdim}(\R^{\adim})$ with $\|V\|(\R^{\adim}) < \infty$ we set $\mathbf{F}(V) =
    \textint{}{} F(S) \ud V(x,S)$, which is~\cite[(1.3)]{DePhilippisGennaioliPigatiRindler2026}.
\end{miniremark}

\begin{definition}
    \label{def:BF}
    For $S \in \pgrass{\adim}{\vdim}$ we define $B_F(S) \in \End{\R^{\adim}}$ by
    \begin{displaymath}
        B_F(S) \bullet L = F(S) \, (S \bullet L) + \uD F(S) \bigl( S^\perp \circ L \circ S +
        \transpose{(S^\perp \circ L \circ S)} \bigr)
    \end{displaymath}
    for $L \in \End{\R^{\adim}}$;
    the argument of $\uD F(S)$ belongs to $\Tan(\pgrass{\adim}{\vdim},S)$
    by~\cite[Lemma~A.1]{DePhilippis2018}. This is~\cite[(A.5)]{DePhilippis2018} for an~integrand
    not depending on the point. We set $K_F = \conv B_F[\pgrass{\adim}{\vdim}]$, and say that $F$
    satisfies \emph{(AC1)} if $\dim \im A \ge \vdim$ for every $A \in K_F$;
    see~\cite[section~1.2]{DePhilippisGennaioliPigatiRindler2026}. If $F$ satisfies (AC1), we set
    \begin{displaymath}
        \lambda(F) = \inf \{ \sv{\vdim}(A) : A \in K_F \} \,, \qquad
        \Lambda(F) = \sup \{ |A| : A \in K_F \} \,.
    \end{displaymath}
\end{definition}

\begin{lemma}
    \label{lem:first_variation}
    Suppose $V \in \Var{\vdim}(\R^{\adim})$, $\|V\|(\R^{\adim}) < \infty$, $g \in
    \VF(\R^{\adim})$, and $\phi_t(x) = x + t g(x)$ for $t \in \R$ and $x \in \R^{\adim}$. Then
    there is $0 < t_0 < \infty$ such that $\phi_t$ is a~diffeomorphism of~$\R^{\adim}$ onto
    itself for $|t| < t_0$, the derivative of $[t \mapsto \mathbf{F}(\phi_{t\#} V)]$ at $t = 0$
    exists, and
    \begin{displaymath}
        \delta_F V(g) = \frac{\ud}{\ud t} \mathbf{F}(\phi_{t\#} V) \Big|_{t=0} = \textint{}{}
        B_F(S) \bullet \uD g(x) \ud V(x,S) \,.
    \end{displaymath}
\end{lemma}

\begin{proof}
    The first assertion is stated in the proof of~\cite[Lemma~A.2]{DePhilippis2018}.
    By~\cite[3.2(1)]{Allard1972}, for $|t| < t_0$,
    \begin{displaymath}
        \mathbf{F}(\phi_{t\#} V) = \textint{}{} F(S_t(x,S)) \, \|\extpower{\vdim}(\uD \phi_t(x)
        \circ S)\| \ud V(x,S) \,,
    \end{displaymath}
    where $S_t(x,S) \in \pgrass{\adim}{\vdim}$ is the orthogonal projection onto $\uD \phi_t(x)[\im
    S]$. Since $\uD \phi_t(x) = \id{\R^{\adim}} + t \uD g(x)$, we have
    \begin{displaymath}
        \frac{\ud}{\ud t} \|\extpower{\vdim}(\uD \phi_t(x) \circ S)\| \Big|_{t=0} = \uD g(x) \bullet
        S \,, \qquad
        \frac{\ud}{\ud t} S_t(x,S) \Big|_{t=0} = S^\perp \circ \uD g(x) \circ S + \transpose{(S^\perp
        \circ \uD g(x) \circ S)}
    \end{displaymath}
    by~\cite[2.4(1)]{Allard1972} and~\cite[Lemma~A.1]{DePhilippis2018}. The integrand equals
    $F(S)$ for $x \notin \spt g$, and its derivative in~$t$ is bounded on $\{ t : |t| \le t_0/2
    \} \times \R^{\adim} \times \pgrass{\adim}{\vdim}$. Differentiation under the integral sign,
    justified by the mean value theorem and~\cite[2.4.9]{Federer1969}, and~\ref{def:BF} give the
    conclusion.
\end{proof}

\begin{lemma}
    \label{lem:KF}
    The set $K_F$ is compact and convex. If $F$ satisfies (AC1), then $0 < \lambda(F) \le
    \Lambda(F) < \infty$, $\sv{\vdim}(A) \ge \lambda(F)$ and $|A| \le \Lambda(F)$ for $A \in K_F$.
\end{lemma}

\begin{proof}
    The map~$B_F$ is continuous, because $F$ is of class~$1$; hence $B_F[\pgrass{\adim}{\vdim}]$
    is compact, and $K_F$ is compact by~\ref{lem:hull}. Suppose $F$ satisfies (AC1).
    By~\ref{lem:sv}\ref{i:sv:rank}, $\sv{\vdim}(A) > 0$ for $A \in K_F$. Since $\sv{\vdim}$ is
    continuous by~\ref{lem:sv}\ref{i:sv:lip}, it attains a~positive minimum on~$K_F$; hence
    $\lambda(F) > 0$. Since $\sv{\vdim}(A) \le \|A\| \le |A|$, we have $\lambda(F) \le
    \Lambda(F)$, and $\Lambda(F) < \infty$ by compactness.
\end{proof}

\begin{theorem}
    \label{thm:MS}
    Suppose $F$ satisfies (AC1). There is $0 < C < \infty$, depending only on $\adim$,
    $\lambda(F)$ and $\Lambda(F)$, with the following property. Suppose $V \in
    \RVar{\vdim}(\R^{\adim})$, $\|V\|(\R^{\adim}) < \infty$, $\density^{\vdim}(\|V\|,x) \ge 1$ for
    $\|V\|$~almost all~$x$, the linear functional $\delta_F V$ on~$\VF(\R^{\adim})$
    of~\ref{lem:first_variation} is representable by integration, $\|\delta_F V\|(\R^{\adim}) <
    \infty$, and $M = \R^{\adim} \cap \{ x : 1 \le \density^{\vdim}(\|V\|,x) < \infty \}$. Then
    \begin{gather}
        \|V\|(\R^{\adim}) \le C \, \HM^{\vdim}(M)^{1/\vdim} \, \|\delta_F V\|(\R^{\adim}) \,,
        \label{eq:MS:MS} \\
        \|V\|(\R^{\adim})^{(\vdim - 1)/\vdim} \le C \, \|\delta_F V\|(\R^{\adim})
        \quad \text{if $\vdim > 1$} \,.
        \label{eq:MS:sobolev}
    \end{gather}
\end{theorem}

\begin{proof}
    Let $C$ be the constant of~\ref{thm:main} for $Y = \R^{\adim}$, $\lambda = \lambda(F)$ and
    $\Lambda = \Lambda(F)$. The functions $[x \mapsto \|V\| \cball{x}{r}]$ are upper
    semicontinuous for $0 < r < \infty$, and $[r \mapsto \|V\| \cball{x}{r}]$ is continuous from
    the right; hence the upper and lower densities of~$\|V\|$ are Borel functions, and $M$ is the
    Borel set of points at which both are equal, finite, and at least~$1$. By~\cite[2.10.19(3)]{Federer1969}, applied to the measure~$\|V\|$, the
    set~$M$, the open set~$\R^{\adim}$ and a~number $0 < \tau < 1$, and by $\HM^{\vdim} \le
    \mathscr{S}^{\vdim}$~\cite[2.10.6]{Federer1969}, we have $\tau \HM^{\vdim}(M) \le
    \|V\|(\R^{\adim})$; letting $\tau \to 1-$,
    \begin{equation}
        \label{eq:MS:HM}
        \HM^{\vdim}(M) \le \|V\|(\R^{\adim}) < \infty \,.
    \end{equation}

    By~\cite[3.5]{Allard1972}, $V = \textsum{i=1}{\infty} c_i \, \var{}(E_i)$ for some positive
    numbers $c_i$ and $\HM^{\vdim}$~measurable sets $E_i$ meeting every compact set in
    an~$(\HM^{\vdim},\vdim)$ rectifiable set. Hence $A = \tbcup_{i=1}^\infty E_i$ is countably
    $(\HM^{\vdim},\vdim)$ rectifiable~\cite[3.2.14]{Federer1969} and $\|V\|(\R^{\adim} \without A)
    = 0$; thus $\|V\| \restrict A = \|V\|$. By~\cite[2.10.19(4)]{Federer1969} applied to $\|V\|$
    and the $\|V\|$~measurable set~$A$, and by $\HM^{\vdim} \le \mathscr{S}^{\vdim}$, we obtain
    $\density^{\vdim}(\|V\|,x) = 0$ for $\HM^{\vdim}$~almost all $x \in \R^{\adim} \without A$;
    hence $\HM^{\vdim}(M \without
    A) = 0$, and $M$ is countably $(\HM^{\vdim},\vdim)$ rectifiable.

    By~\cite[3.5(1)]{Allard1972}, for $\|V\|$~almost all~$x$ the density
    $\density^{\vdim}(\|V\|,x)$ exists and is finite, $\Tan^{\vdim}(\|V\|,x) \in
    \pgrass{\adim}{\vdim}$, $\|V\| = \HM^{\vdim} \restrict \density^{\vdim}(\|V\|,\cdot)$, and
    \begin{equation}
        \label{eq:MS:V}
        V(\alpha) = \textint{}{} \alpha(x,\Tan^{\vdim}(\|V\|,x)) \, \density^{\vdim}(\|V\|,x) \ud
        \HM^{\vdim} x
        \quad \text{for $\alpha \in \ccspace{\R^{\adim} \times \pgrass{\adim}{\vdim}}$} \,.
    \end{equation}
    By hypothesis and~\cite[3.5(1)]{Allard1972}, $\|V\|(\R^{\adim} \without M) = 0$. Let $\theta(x) = \density^{\vdim}(\|V\|,x)$
    for $x \in M$; then $\|V\| = (\HM^{\vdim} \restrict M) \restrict \theta$, and $\theta$ is
    $\HM^{\vdim} \restrict M$ summable. Fix $S_0 \in \pgrass{\adim}{\vdim}$, and let $X(x) =
    B_F(\Tan^{\vdim}(\|V\|,x))$ for those $x \in M$ at which $\Tan^{\vdim}(\|V\|,x) \in
    \pgrass{\adim}{\vdim}$, and $X(x) = B_F(S_0)$ for the other $x \in M$. Then $X : M \to K_F$ is
    $\HM^{\vdim} \restrict M$ measurable by~\cite[3.5(1)]{Allard1972}. Let $\sigma(\Phi) =
    \textint{}{} X \bullet \Phi \ud \|V\|$ for $\Phi \in
    \ccspace{\R^{\adim},\End{\R^{\adim}}}$, and $\Upsilon = - \delta_F V$. For $g \in
    \VF(\R^{\adim})$ the function $[(x,S) \mapsto B_F(S) \bullet \uD g(x)]$ belongs to
    $\ccspace{\R^{\adim} \times \pgrass{\adim}{\vdim}}$; hence, by~\ref{lem:first_variation}
    and~\eqref{eq:MS:V},
    \begin{displaymath}
        \Upsilon(g) = - \textint{}{} B_F(S) \bullet \uD g(x) \ud V(x,S) = - \sigma(\uD g) \,.
    \end{displaymath}
    Moreover, $\|\Upsilon\| = \|\delta_F V\|$. By~\ref{lem:KF}, the data $Y = \R^{\adim}$, $K =
    K_F$, $\lambda(F)$, $\Lambda(F)$, $M$, $\theta$, $X$, $\mu = \|V\|$, $\sigma$, $\Upsilon$
    satisfy~\ref{mr:hypotheses_main}. Hence~\ref{thm:main} gives~\eqref{eq:MS:MS}.
    Suppose $\vdim > 1$. If $\|V\|(\R^{\adim}) > 0$, combining~\eqref{eq:MS:MS}
    with~\eqref{eq:MS:HM} and dividing by $\|V\|(\R^{\adim})^{1/\vdim}$ gives~\eqref{eq:MS:sobolev};
    otherwise~\eqref{eq:MS:sobolev} is trivial.
\end{proof}

\section{Dependencies and formalisation notes}
\label{sec:summary}

\begin{miniremark}
    \label{mr:sec_summary}
    This section lists the dependencies of all items in Table~\ref{tab:dependencies} and records
    which parts of the deduction are available in existing libraries.
\end{miniremark}

\begin{table}[ht]
    \centering
    \small
    \begin{tabular}{p{0.06\textwidth}p{0.3\textwidth}p{0.54\textwidth}}
        item & source & uses \\ \hline
        \ref{lem:sv} & new & \cite[1.7.3, 1.7.9]{Federer1969} \\
        \ref{lem:spectral} & new & \cite[1.7.3, 1.7.9]{Federer1969} \\
        \ref{lem:trace} & new & \cite[1.7.3, 1.7.10]{Federer1969} \\
        \ref{lem:hull} & new (Carath\'eodory) & --- \\
        \ref{lem:barycentre} & --- & \cite[2.4.12, 2.4.19]{Federer1969} \\
        \ref{lem:kernel} & new & \cite[2.6.2]{Federer1969} \\
        \ref{lem:heat_extension} & new & \cite[2.4.9, 2.4.12, 2.6.2]{Federer1969},
        \ref{lem:kernel} \\
        \ref{lem:classes} & new & \cite[2.4.9, 2.6.2, 3.1.11]{Federer1969} \\
        \ref{lem:semiconvex} & DGPR (3.16) & \cite[2.4.14]{Federer1969}, \ref{lem:kernel},
        \ref{lem:heat_extension} \\
        \ref{lem:ibp} & DGPR (3.17) & \cite[2.4.9, 2.9.9, 3.1.6, 4.1.2]{Federer1969},
        \ref{lem:classes} \\
        \ref{lem:identities} & DGPR (3.1)--(3.4), GR (3.9) & \cite[2.5.14, 4.1.5]{Federer1969},
        \ref{lem:sv}, \ref{lem:barycentre}, \ref{lem:heat_extension} \\
        \ref{lem:regularity} & --- & \cite[3.1.11]{Federer1969}, \ref{lem:heat_extension},
        \ref{lem:classes}, \ref{lem:identities} \\
        \ref{lem:heat_calculus} & DGPR (3.6) & \cite[3.1.11]{Federer1969}, \ref{lem:identities}
        \\
        \ref{lem:DGPR31} & DGPR Lemma~3.1 & \cite[2.4.7, 2.4.19]{Federer1969},
        \ref{lem:classes}, \ref{lem:identities}--\ref{lem:heat_calculus} \\
        \ref{lem:DGPR32} & DGPR Lemma~3.2 & \ref{lem:sv}, \ref{lem:spectral}, \ref{lem:trace},
        \ref{lem:classes}--\ref{lem:ibp}, \ref{lem:identities}, \ref{lem:regularity},
        \ref{lem:DGPR31} \\
        \ref{lem:entropy} & DGPR section~4 & \ref{lem:kernel}, \ref{lem:classes},
        \ref{lem:identities}, \ref{lem:regularity} \\
        \ref{lem:upper_entropy} & DGPR (4.1) & \ref{lem:DGPR32}, \ref{lem:entropy} \\
        \ref{lem:tube} & DGPR proof of (4.2) & \cite[2.6.2, 2.7.16, 3.2.37, 3.2.39]{Federer1969} \\
        \ref{lem:lower_entropy} & DGPR (4.2) & \cite[2.4.19, 2.6.2]{Federer1969},
        \ref{lem:kernel}, \ref{lem:identities}, \ref{lem:regularity}, \ref{lem:tube} \\
        \ref{lem:inner} & new & \cite[2.2.2, 2.10.43, 3.2.14]{Federer1969} \\
        \ref{lem:compact_case} & DGPR proof of Theorem~1.1 & \ref{lem:upper_entropy},
        \ref{lem:lower_entropy}, \ref{lem:inner} \\
        \ref{thm:main} & DGPR Theorem~1.1 for rectifiable~$M$ & \ref{mr:functionals}, \ref{lem:compact_case} \\
        \ref{lem:first_variation} & DDG Lemma~A.2 & \cite[2.4(1), 3.2(1)]{Allard1972},
        \cite[Lemma~A.1]{DePhilippis2018}, \cite[2.4.9]{Federer1969} \\
        \ref{lem:KF} & --- & \ref{lem:sv}, \ref{lem:hull} \\
        \ref{thm:MS} & DGPR Theorem~1.6 & \cite[3.5]{Allard1972}, \cite[2.10.6, 2.10.19,
        3.2.14]{Federer1969}, \ref{thm:main}, \ref{lem:first_variation}, \ref{lem:KF}
    \end{tabular}
    \caption{Dependencies. DGPR is~\cite{DePhilippisGennaioliPigatiRindler2026}, GR
    is~\cite{Gennaioli2026}, DDG is~\cite{DePhilippis2018}; ``new'' marks results which are in
    none of the three sources.}
    \label{tab:dependencies}
\end{table}

\begin{remark}
    \label{rem:rocq}
    For a~formalisation in Rocq on top of a~formalisation of Federer's book, the items to be
    added are those marked ``new'' or attributed to DGPR in Table~\ref{tab:dependencies}, the
    formalisation of~\cite[2.3--4.3]{Allard1972} cited in section~\ref{sec:varifolds}, and
    \cite[Lemmas~A.1 and~A.2]{DePhilippis2018}. All other inputs are numbered items of
    Federer's book. The heat flow layer of section~\ref{sec:heat} is self-contained: it requires
    only integration with respect to Borel measures, Fubini's theorem, differentiation under the
    integral sign, and the one-dimensional Gaussian integrals. Positivity of the heat kernel is
    used in~\ref{lem:heat_extension}\ref{i:he:gauss}, \ref{lem:semiconvex}
    and~\ref{lem:regularity}; hence the heat extension must be formalised for measures, not
    merely for functions in~$\Lp{p}$.
\end{remark}

\begin{remark}
    \label{rem:lean}
    For a~formalisation in Lean, Mathlib provides Lebesgue and Hausdorff measures, Fubini's
    theorem, Fr\'echet derivatives, Jensen's inequality, Rademacher's theorem, Lebesgue's
    differentiation theorem, Gaussian integrals and Carath\'eodory's theorem on convex hulls; the
    exact declaration names should be checked against the current Mathlib. The library Tau~Ceti
    contains countably rectifiable sets (\url{TauCeti/Geometry/MeasureTheory/Rectifiable/Basic.lean}),
    varifolds and their weights (\url{TauCeti/Geometry/MeasureTheory/Varifold/Euclidean.lean}),
    the heat semigroup on~$\Lp{p}$ (\url{TauCeti/Analysis/PDE/HeatKernel/Semigroup.lean}),
    an~entropy functional (\url{TauCeti/MeasureTheory/Integral/Entropy.lean}), and limits of
    volumes of thickenings of compact sets
    (\url{TauCeti/Analysis/BoundedVariation/MinkowskiContent.lean}). The following inputs were
    not found there and have to be formalised: distributions representable by integration
    (\ref{mr:functionals}), the Minkowski content of closed rectifiable
    sets~\cite[3.2.39]{Federer1969}, heat extensions of measures (\ref{lem:heat_extension}), and
    Allard's image varifolds and rectifiable varifolds~\cite[3.2, 3.5]{Allard1972}.
\end{remark}

\bibliographystyle{halpha}
\bibliography{blueprint}

\end{document}
